% Canonical manuscript for the website and PDF. Edit this file.

% TeX fragment: compile through the edition builder (make all).

\title{formlære}

\emph{Ei traktatform om form for design og estetiske fag}\newline{}Iver Raknes Finne · Arkitektur- og designhøgskolen i Oslo · 2026

\section{Føreord}

Designfaget har eit stort arkiv av metodar, kritikk, historie og arbeid gjennom praksis. Likevel står det grunnleggjande spørsmålet ope: Kva er ei form, kva kunne ha stått i staden for henne, og kvifor vart nettopp denne forma realisert? Det finst forskingsprogram i design, men ikkje eitt felles, kumulativt og prøvbart program som bind desse spørsmåla saman. Eit fag er ikkje grunnlagt berre fordi det kan skildre korleis utøvarane arbeider. Det må òg kunne seie kva som ville ha motbevist forklaringa deira.\kjelder{153,154,155,162,169} Denne boka er eit forslag til eit slikt program.

To sirkelargument står i vegen. Det første seier at forma følgjer funksjonen. Men ei delt oppgåve kan ikkje åleine forklare skilnadene mellom formene som løyser henne. Dersom «funksjon» først vert definert etter at forma er vald, seier læresetninga berre at valet svara til valet. For å forklare må ein i staden vise kva alternativ som var tilgjengelege, og kva vilkår som skilde dei.\kjelder{125,15,77,78} Den same slutningsfeilen oppstår når kategoriane ein vil stadfeste, alt er bygde inn i målemetoden. Historisk raseantropologi viser kor alvorlege følgjer ein slik logisk sirkel kan få; samanlikninga gjeld slutningsmåten, ikkje ei likestilling av fag eller skade.\kjelder{163,164}

Det andre argumentet gjer prosessen til bevis. Ei skisse viser at nokon skisserte. Ein prototype viser kva som vart bygd. Ingen av delane viser, utan ei uavhengig prøve, kva som kunne ha hendt i staden, eller kvifor forma verka. Simon, Rittel og Webber, Schön og Cross opna avgjerande spørsmål om handling, problem, refleksjon og fagleg kunnskap.\kjelder{119,160,161,162} Galle har undersøkt grunnlaget for designteori, Redström korleis teori vert gjord, og forsking gjennom design korleis innsikt kan vekse frå artefakt og praksis.\kjelder{126,127,128,138,154,155} Systemorientert design har gjeve reiskapar til å sjå samanhengar og verknader på tvers av skalaer.\kjelder{129,130,146} Bidraga er verkelege. Men ingen metode får forklaringskraft av namnet sitt. Motdømet må kunne kome utanfrå metoden.

Dette er ikkje berre eit vitskapsteoretisk spørsmål. Eit grensesnitt kan vere glimrande utforma for å halde eit barn fast. Ei teneste kan oppfylle bestillinga ved å hente ut data og gjere utmelding vanskeleg. Ein gjenstand kan vere billeg å kjøpe og dyr å bryte ned.\kjelder{132,134,142} Når skaden ligg i det bestilte, er «krava oppfylte» ingen dom over forma. Den som lever med følgjene, har ikkje nødvendigvis valt krava. Eit fag som berre kan stadfeste bestillinga, kan ikkje grunngje motstanden mot henne.

formlære tek utgangspunkt i eit enklare spørsmål: Kva skilnader er ordna, kva alternativ var moglege, og kva gjorde utslaget? Svaret er ein substratuavhengig teori om formdanning gjennom samspelet mellom generative grammatikkar, kunnskapsutvikling og agens på fleire skalaer. Form er ei ordning av skilnader. Formgrammatikken gjev stega som kan generere alternativ. Landskapet gjev vilkåra som vektar dei. Agentane gjer avvik, mål og avgrensa kapasitet til handling. Kunnskapsstanden skil det kjende frå det tenkte og det som enno ikkje kan tenkjast. Materialet avgrensar kva som kan realiserast. Desse skilja vert haldne saman utan å blandast saman.

Arbeidet byggjer på, men er ikkje identisk med, teknikk og tektonikk hos Semper og Pye;\kjelder{1,26} morforom hos Thompson og Raup;\kjelder{2,10} generative reglar hos Stiny og Gips;\kjelder{11,18} landskap hos Wright og Waddington;\kjelder{3,8} tilbakekopling hos Wiener og Ashby;\kjelder{4,5,7} skiljet mellom konsept og kunnskap hos Hatchuel og Weil;\kjelder{16,17} og statistisk inferens og fleirskala agens hos Jaynes, Friston, Fields og Levin.\kjelder{19,21,22,83,86,113} Arkitekturen vert samla som \((\mathbb{M},\Pi,SG,\Lambda,\{A_i\})\). Eit felles namn på desse ledda er enno ikkje eit bevis. Poenget er at dei no kan fastsetjast, skiljast og prøvast.

Grensetilfella viser kva som står att når eitt ledd vert halde fast eller teke bort. Dei gjer ikkje heile formgrammatikken, C–K-teorien, teknisk optimering eller Alexander si formsyntese til særtilfelle av éi oppdaga naturlov.\kjelder{119,120} Stil kan både vere eit statistisk mønster og eit kulturelt grensesnitt, men dei to forklaringane er ikkje identiske.\kjelder{9,53,80} Eit skissesteg kan trekkjast tilbake; ein materiell overgang kan ikkje alltid det. Difor må likevekt i søket skiljast frå historie i stoffet.\kjelder{14}

I ein agent \(A=(m,\pi,G,\varepsilon)\) må me skilje kva han føreseier, kva han bryr seg om, kva han kan halde, og kva som faktisk hender.\kjelder{21,23,68} Meir prioritering kan minske det \emph{estimerte} avviket og samstundes auke den verkelege skaden dersom estimatet peikar feil (6.673–6.674). Å kjenne til eit omsyn er ikkje å halde det. Å ønskje å halde det er heller ikkje å kunne. Blindskap, likegyldigheit og avmakt får difor kvar si prøve. Omhug er at avvik set agenten i arbeid. Omsorg er den relasjonelle utvidinga av målrommet: avviket hos ein annan vert teke inn i det agenten sjølv arbeider for å minske (5.28). Kvar agent som realiserer former, er ein formgjevar; ein designar er ein formgjevar med omsorg for dei forma verkar på (6.01–6.02). Når omsorga vert bygd ut på tvers av skalaer, gjev den same matematiske arkitekturen som forklarer form, intelligens og samarbeid, òg eit strukturelt grunnlag for etisk formgjeving (6.7, 7.26).

formlære er eit forskingsprogram, ikkje eit ferdig system (7.17). Den harde kjernen er påstandane arkitekturen kviler på. Kjernelikninga i 5.7 er ein hypotese om lokale steg; fordelinga i 6.6 er ein valmodell under uttrykkelege vilkår om grunnfordeling, estimat og normalisering.\kjelder{110,113,121,122} Begge er hjelpehypotesar og kan fellast utan at kjernen fell. Prøva av læra er ikkje at ho kan forklare den forma som alt finst. Ho må føreseie nye val, nye sviktar og verknader av endra vilkår betre enn ein oppgjeven alternativ modell. Når ho feilar, skal me kunne finne om feilen låg i stega, vurderinga, prioriteringa, kapasiteten eller følgjene. Den andre sitt svar må kunne rette både modellen og målet. Slik vert etikken ikkje ei høgtideleg avslutning på eit arbeid som alt har gjort skade, men eit krav til kva arbeidet må halde (7.26).\notat{intro-1}

Eit fundament skal ikkje gje siste ordet. Det skal gjere det mogleg å finne nøyaktig kvar eit ord ikkje held. Forma stoggar. Verknaden held fram.

\emph{Oslo, 2026}


\section{Lesekode}

d \(=\) definisjon · a \(=\) postulat eller hypotese · t \(=\) følgje under oppgjevne vilkår · o \(=\) observasjon · i \(=\) illustrasjon. Desimalnummera viser logisk vekt: 2.1 kommenterer 2, og 2.11 kommenterer 2.1. Nummer med null, som 2.01, er utdjupingar som står før 2.1. Dei sju heiltalsproposisjonane kan lesast som eitt argument. Definisjonar fastset språkbruken. Eit postulat vel ein ståstad; ein hypotese seier kva som ville ha felt han. Ei følgje står og fell med premissane sine. Illustrasjonar ber inga slutning.

\begin{formproposition}{1}{d}
Form er ei ordning av skilnader.\kjelder{1}
\end{formproposition}

\begin{formproposition}{1.01}{d}
Å skilje noko ut er å halde noko anna utanfor. Det som då står mellom dei, er ein skilnad.
\end{formproposition}

\begin{formproposition}{1.011}{t}
Ein skilnad har difor to sider. Han gjev ikkje to uavhengige verder, for kvar side er det ho er, berre mot den andre.
\end{formproposition}

\begin{formproposition}{1.0111}{i}
Kanten på koppen skil innsida frå utsida. Munnen møter begge ved den same kanten.
\end{formproposition}

\begin{formproposition}{1.012}{o}
Ein skilnad kan verke før han får eit namn. (Membranen treng ikkje ordet for å halde noko inne.)
\end{formproposition}

\begin{formproposition}{1.0121}{t}
Å oppdage ein skilnad er difor ikkje å lage han. Ei teikning kan gjere begge delar.
\end{formproposition}

\begin{formproposition}{1.02}{d}
Skilnader som står i bestemte samanhengar med kvarandre, er ei ordning: kva som er skilt frå kva, og korleis det skilde heng saman.
\end{formproposition}

\begin{formproposition}{1.021}{t}
Delane i ei ordning er ikkje alltid ferdige før ordninga.

Kvar side av ein skilnad er det ho er mot den andre (1.011). Det som står i ei ordning, kan difor få si bestemming gjennom ordninga.
\end{formproposition}

\begin{formproposition}{1.0211}{i}
Eit sete bind saman ei bereflate, ein kropp og ei mogleg kvile. Flata kan finnast utan kroppen. Setet kan ikkje forklarast utan kroppen.
\end{formproposition}

\begin{formproposition}{1.1}{a}
Formverda er alt som er tilfelle: totaliteten av realiserte konfigurasjonar, ikkje av ting.
\end{formproposition}

\begin{formproposition}{1.11}{d}
Ein konfigurasjon er ei bestemt ordning av bestemte delar (1.02). Han er realisert når han er tilfelle.
\end{formproposition}

\begin{formproposition}{1.111}{i}
Same treet kan vere stol eller ved.
\end{formproposition}

\begin{formproposition}{1.12}{t}
Kvar realisert konfigurasjon er bestemt (1.11). Anten er han tilfelle, eller han er det ikkje.

Ingen konfigurasjon er delvis realisert.
\end{formproposition}

\begin{formproposition}{1.13}{d}
Å forklare kvifor ein bestemt konfigurasjon vart realisert, er å vise kva vilkår og overgangar som førte til han, og kva dei gjorde med alternativa.
\end{formproposition}

\begin{formproposition}{1.131}{t}
Forklaringa er difor òg ei forklaring av utvalet. Ho må skilje mellom alternativ som var umoglege, utilgjengelege, uprøvde og forkasta. Desse fråværa har ulike årsaker.
\end{formproposition}

\begin{formproposition}{1.2}{d}
Ein konfigurasjonsdel som ikkje vert delt vidare på det valde analysenivået, er eit objekt.

Same objekt kan stå i fleire konfigurasjonar.
\end{formproposition}

\begin{formproposition}{1.201}{t}
Inndelinga i objekt er eit val av analysenivå.
\end{formproposition}

\begin{formproposition}{1.2011}{i}
Stolbeinet er éin del for snikkaren som monterer. For den som granskar brotet, er fibrane delar. Nivået må følgje det som skal forklarast.
\end{formproposition}

\begin{formproposition}{1.202}{a}
Kvart objekt er ein konfigurasjon på eit lågare nivå. Under kvart valt nivå held strukturen fram.
\end{formproposition}

\begin{formproposition}{1.2021}{t}
Ein konfigurasjon er ein samanheng (1.11). Det analysen finn, er difor relasjonar.
\end{formproposition}

\begin{formproposition}{1.2022}{a}
Ein relasjon krev ledd. At ledda sjølve kan analyserast som relasjonar, gjer ikkje spørsmålet om kva relasjonane er i, til endå ei oppdeling.
\end{formproposition}

\begin{formproposition}{1.203}{d}
Substans er den indre naturen til det som står i relasjonane. Han er ikkje eit ekstra ledd i den fysiske skildringa, men det ledda er, teke uavhengig av korleis dei er ordna.
\end{formproposition}

\begin{formproposition}{1.2031}{t}
Substansen er ikkje det siste objektet analysen finn. Analysen kan halde fram gjennom nivå utan ende, medan spørsmålet om substans gjeld kva som er ordna på alle nivåa.
\end{formproposition}

\begin{formproposition}{1.204}{t}
Objekta vert haldne faste på det valde nivået; konfigurasjonane skiftar. Endrar eit fysisk stykke ein eigenskap analysen held fast, vert endringa skildra som ein annan objekttilstand.
\end{formproposition}

\begin{formproposition}{1.21}{d}
Måten objekta heng saman på i konfigurasjonen, er strukturen til konfigurasjonen.
\end{formproposition}

\begin{formproposition}{1.211}{i}
Sambandet mellom setet og beinet i ein stol er ikkje ein del til. Delane heng i kvarandre direkte, som lenkene i ei kjede.
\end{formproposition}

\begin{formproposition}{1.212}{d}
Forma til ein konfigurasjon er strukturen hans, teken uavhengig av dei eigenskapane ved objekta analysen ser bort frå.

Abstraksjonen må oppgjevast. Same silhuett er ikkje nødvendigvis same bereform.
\end{formproposition}

\begin{formproposition}{1.2121}{t}
To konfigurasjonar med same objekt, men ulik struktur, har ulike former.
\end{formproposition}

\begin{formproposition}{1.2122}{t}
To konfigurasjonar med ulike objekt, men same struktur, deler form i dei omsyna abstraksjonen held fast ved.
\end{formproposition}

\begin{formproposition}{1.2123}{t}
Forma kan difor gå att i andre objekt. Verknadene går att berre i den grad dei òg kviler på dei relasjonane som er bevarte.
\end{formproposition}

\begin{formproposition}{1.2124}{t}
Ei form utan alternativ under gjevne vilkår er nødvendig relativt til vilkåra. Kvifor nettopp desse vilkåra gjeld, kan framleis krevje ei forklaring.
\end{formproposition}

\begin{formproposition}{1.22}{d}
Moglegheita eit objekt på eit valt analysenivå (1.201) har for å inngå i konfigurasjonar på dette nivået, er valensen til objektet.
\end{formproposition}

\begin{formproposition}{1.221}{t}
Når objektdelinga, relasjonane og operasjonane er fastsette, er dei kombinatoriske formene på dette nivået gjevne i og med valensane til objekta. At rommet er fastsett, seier ikkje at alle posisjonane kan reknast opp.
\end{formproposition}

\begin{formproposition}{1.222}{a}
Det moglege er ikkje laga, men gjeve.

Inga objektdeling gjev eit fullstendig kart på førehand.
\end{formproposition}

\begin{formproposition}{1.2221}{t}
Ei moglegheit i feltet kan ikkje finnast opp, berre avdekkjast. Det som vert laga, er realiseringa.
\end{formproposition}

\begin{formproposition}{1.2222}{t}
Det vidare moglege ligg fast; valet av nivå, objekt, aksar og operasjonar avgrensar kartet over det.

Eit nivåskifte kan utvide det operative formrommet utan å avgjere kva som er ontologisk mogleg.
\end{formproposition}

\begin{formproposition}{1.223}{t}
Valensen høyrer til objektet, og objektet står fast på det valde nivået (1.204). Ei ny kombinatorisk moglegheit for dette objektet kjem ikkje til på same nivå under uendra vilkår.

Ei inndeling på tvers av objekta (1.611) eller eit skifte av analysenivå (1.201) kan avdekkje moglegheiter i feltet som den førre objektdelinga skjulte. Nye relasjonar eller operasjonar kan òg endre det operative formrommet.
\end{formproposition}

\begin{formproposition}{1.2231}{t}
Realisering, tilgang og kunnskap kan endre seg utan at det moglege endrar seg.
\end{formproposition}

\begin{formproposition}{1.23}{a}
Substansen i formverda er oppleving. Inga fysisk måling kan avgjere det.
\end{formproposition}

\begin{formproposition}{1.231}{d}
Ei måling er ein relasjon mellom det som vert målt og det som måler.
\end{formproposition}

\begin{formproposition}{1.2311}{t}
Det ei måling gjev, er difor ein relasjon, aldri leddet sjølv.
\end{formproposition}

\begin{formproposition}{1.232}{o}
Fysikken seier alt han seier i storleikar, og kvar storleik er eit utfall av måling eller ein relasjon mellom slike utfall.
\end{formproposition}

\begin{formproposition}{1.2321}{t}
Fysikken seier difor kva det målte gjer, ikkje kva det er.
\end{formproposition}

\begin{formproposition}{1.233}{d}
Det som er kjent ved å vere det, ikkje gjennom ein målerelasjon (1.231), er oppleving.

Dette skil tilgangsmåtar. Det seier ikkje at oppleving manglar indre skilnader eller relasjonar.
\end{formproposition}

\begin{formproposition}{1.2331}{o}
Kvar måling er gjeven som oppleving. Det fysikken veit, veit han gjennom opplevinga.
\end{formproposition}

\begin{formproposition}{1.2332}{d}
Oppleving av opplevinga, opplevinga framstilt på nytt i seg sjølv, er medvit.
\end{formproposition}

\begin{formproposition}{1.2333}{a}
Oppleving føreset ikkje medvit.
\end{formproposition}

\begin{formproposition}{1.234}{a}
Når ein kjend natur strekkjer til som indre natur for dei relasjonane som skal forklarast, legg denne læra ingen annan substans til.

Det avgjer ikkje kva som finst.
\end{formproposition}

\begin{formproposition}{1.2341}{a}
Den relasjonelle strukturen i den fysiske skildringa kan vere ordning i oppleving. Opplevinga strekkjer til som indre natur utan at dei målte relasjonane må endrast.
\end{formproposition}

\begin{formproposition}{1.2342}{t}
Når ei måling berre gjev relasjonar (1.2311), avgjer ho ikkje kva ledda i relasjonane er. Oppleving er ein natur som er kjend utan måling (1.233). Dersom ho kan vere den indre naturen utan at den fysiske skildringa vert endra (1.2341), krev modellen ikkje ein annan, ukjend substans. Minimalitetskravet i 1.234 vel då oppleving. Dette er grunngjevinga for postulatet i 1.23, ikkje eit bevis for det.

Valet er ikkje tvinga fram av måledata. Andre ontologiar kan gje dei same relasjonane. Dei operasjonelle modellane om grammatikk, overgangar, kapasitet og svikt må stå eller falle på sine eigne, prøvbare premissar.
\end{formproposition}

\begin{formproposition}{1.2343}{t}
Oppleving er då ikkje ein eigenskap som eit elles ikkje-opplevande stoff må produsere. Det som skal forklarast vidare, er korleis oppleving er ordna og avgrensa.
\end{formproposition}

\begin{formproposition}{1.235}{d}
Når ei avgrensing skil eit indre frå eit ytre (1.011, 5.02), er det å kjenne henne innanfrå å vere henne, og å kjenne henne utanfrå å kjenne relasjonane ho står i. Tilgangane gjeld same verkelegheit, ikkje to stoff som først må koplast saman.
\end{formproposition}

\begin{formproposition}{1.2351}{t}
Fysikken skildrar relasjonane materien står i (1.2311, 1.2321). Materie er difor oppleving sett utanfrå (1.23). I eitt opplevingsfelt (1.24) er det materielle biletet korleis eit avgrensa mønster tek seg ut for eit anna (1.241).
\end{formproposition}

\begin{formproposition}{1.2352}{t}
Materien er ikkje ein annan ting enn opplevinga, og opplevinga er ikkje ein eigenskap ved materien. Biletet og det det er bilete av, er av same substans.
\end{formproposition}

\begin{formproposition}{1.2353}{t}
Det som vert kjent utanfrå, er relasjonar. Forma er struktur (1.212), og struktur er relasjonar (1.21).

Forma, slik ho vert målt og skildra, er difor korleis ein del av substansen tek seg ut for ein annan.
\end{formproposition}

\begin{formproposition}{1.2354}{t}
At materien er oppleving, tyder ikkje at han er min.
\end{formproposition}

\begin{formproposition}{1.24}{a}
Opplevinga er eitt felt.

Det som vert skilt ut i feltet, er mønster i det, ikkje bitar av det.
\end{formproposition}

\begin{formproposition}{1.2401}{t}
Alle mønstera høyrer til same felt. Kva mønster som kan påverke kvarandre, vert avgjort av relasjonane i det.
\end{formproposition}

\begin{formproposition}{1.241}{a}
Det som er mange frå utsida, kan vere avgrensingar i éin grunn. Der ei avgrensing held seg sjølv (5.021, 5.024), har ho tilgang innanfrå til si eiga avgrensing; andre vert kjende gjennom relasjonane.

Dette tilskriv verken feltet som heilskap ei utside (1.246) eller kvart utskilt mønster eit eige sjølv (5.023).
\end{formproposition}

\begin{formproposition}{1.242}{t}
Mangfaldet er då eit spørsmål om avgrensing i eitt felt. Organiseringa skil ut kva som høyrer saman innanfrå; ho summerer ikkje ferdige opplevingar.
\end{formproposition}

\begin{formproposition}{1.243}{t}
Feltet står fast; konfigurasjonane skiftar.
\end{formproposition}

\begin{formproposition}{1.244}{t}
I feltet kan det skiljast ut mønster som ikkje når kvarandre innanfrå (1.241).
\end{formproposition}

\begin{formproposition}{1.245}{t}
Måten feltet oppfører seg regelmessig på, er naturlovene. Dei vert ikkje lagde på feltet utanfrå, for utanfor feltet finst det ingenting.
\end{formproposition}

\begin{formproposition}{1.2451}{t}
Lovene legg ingen nye ontologiske moglegheiter til. Dei avgrensar kva som kan realiserast under gjevne materielle vilkår.

Når materiale og operasjonar kjem til, kan tilgangen og det operative formrommet endre seg (1.223).
\end{formproposition}

\begin{formproposition}{1.2452}{d}
Materiell tyder: gjeven ved stoff, energi og dei regelmessige samanhengane som avgrensar realiseringa.
\end{formproposition}

\begin{formproposition}{1.2453}{d}
Det som gjeld feltet og lovene i det, same kva nokon veit, er ontologisk. Det som avheng av kva som er kjent om feltet, og av kven, er epistemisk.
\end{formproposition}

\begin{formproposition}{1.246}{t}
Feltet som heilskap har ingenting utanfor seg (1.245). Det kan difor ikkje kjennast utanfrå; det viser seg aldri som ein ting.
\end{formproposition}

\begin{formproposition}{1.3}{d}
Ei mengd konfigurasjonar som tillèt same operasjonar, er ein klasse.
\end{formproposition}

\begin{formproposition}{1.31}{d}
Mengda av alle moglege konfigurasjonar i ein klasse er formrommet \(𝕄\) til klassen.\kjelder{2,10,20}
\end{formproposition}

\begin{formproposition}{1.311}{t}
Det ontologisk moglege vert oppdaga. Det operative formrommet vert avgrensa ved val av klasse, nivå og operasjonar. Denne avgrensinga kan reviderast.
\end{formproposition}

\begin{formproposition}{1.32}{d}
Kvar konfigurasjon i klassen er eitt eksakt punkt i \(𝕄\), posisjonen hans. Kvar akse i formrommet er ein målbar eigenskap.
\end{formproposition}

\begin{formproposition}{1.321}{t}
Eit val av aksar er ein projeksjon av \(𝕄\). Forma finst uavhengig av den einskilde målinga, slik eit fjell finst uavhengig av kartet.
\end{formproposition}

\begin{formproposition}{1.3211}{t}
Fjellet er ikkje i kartet. Men den som les kartet, er av same substans (1.23) som fjellet.
\end{formproposition}

\begin{formproposition}{1.322}{t}
Det empiriske formrommet er alltid ein slik projeksjon. Ei endeleg skildring kan vere tilstrekkeleg for eitt spørsmål og utilstrekkeleg for eit anna. Det som vert utelate, sluttar ikkje å verke.\kjelder{2}
\end{formproposition}

\begin{formproposition}{1.3221}{t}
Projeksjonen er ikkje vilkårleg. Nokre projeksjonar fangar meir variasjon enn andre.
\end{formproposition}

\begin{formproposition}{1.323}{t}
Formrommet er eit rom av relasjonar (1.2311). Det viser korleis konfigurasjonane i klassen tek seg ut utanfrå, ikkje kva dei er innanfrå.
\end{formproposition}

\begin{formproposition}{1.324}{d}
Eit omsyn \(X_k\) er ein akse med eit kriterium: ein målregion og ein toleranse som avviket vert målt mot (5.11). Aksen kan gjelde forma sjølv eller følgjene av henne under oppgjevne vilkår (1.5). Eit omsyn med vekt i eit val verkar som eit trykk (2.1).
\end{formproposition}

\begin{formproposition}{1.4}{d}
Formrommet har eit historisk skilje mellom busett og ikkje busett, og eit materielt skilje mellom tillate og forbode. Det ikkje busette og tillatne er ope.

Skilja må daterast: ei form som var realisert før, kan vere forboden under vilkåra no.
\end{formproposition}

\begin{formproposition}{1.41}{d}
Dei busette regionane er dei posisjonane som er eller har vore realiserte. Dei utgjer det historiske arkivet.
\end{formproposition}

\begin{formproposition}{1.411}{t}
Kvar ny realisering tek utgangspunkt i ein realisert tilstand. Naboar i rommet er dei posisjonane som kan nåast med eitt tillate steg; dei treng ikkje likne kvarandre.
\end{formproposition}

\begin{formproposition}{1.412}{a}
Ein busett posisjon kan trekkje nye realiseringar til seg gjennom verktøy, vane og forventing. Tyngda veks der desse vert haldne oppe; ho kan minke når forma bind knappe ressursar eller fyller behovet.
\end{formproposition}

\begin{formproposition}{1.413}{d}
Totaliteten av busette posisjonar er formhistoria.
\end{formproposition}

\begin{formproposition}{1.42}{d}
Dei forbodne regionane V er posisjonane som er logisk skildra i rommet, men utelukka av dei materielle vilkåra som gjeld.

Forbodet følgjer ikkje regionen åleine, men vilkåra.
\end{formproposition}

\begin{formproposition}{1.421}{t}
Grensa mellom det forbodne og det opne er ikkje logisk, men materiell, for ho er gjeven ved dei regelmessige samanhengane i 1.2452, brukte på dei materiala og operasjonane som er for handa. (Bøk bøyer seg langs fiberen; granitt gjer det ikkje.)
\end{formproposition}

\begin{formproposition}{1.422}{t}
Eit skifte i kva materiale og operasjonar som er for handa, kan opne ein forboden region. Tilgangen veks når nye vegar fører dit (1.2231, 1.431).
\end{formproposition}

\begin{formproposition}{1.43}{d}
Dei opne regionane er dei ikkje busette posisjonane som er materielt tillatne (1.4).
\end{formproposition}

\begin{formproposition}{1.431}{d}
Ein open region som kan nåast frå det busette utan å gå gjennom det forbodne, er tilgjengeleg.
\end{formproposition}

\begin{formproposition}{1.4311}{t}
Ein open region som berre grensar til forbodne regionar, er utilgjengeleg (1.431) utan materialskifte.
\end{formproposition}

\begin{formproposition}{1.432}{d}
Dei opne posisjonane som grensar til det busette, er det nærliggande moglege.\kjelder{13}
\end{formproposition}

\begin{formproposition}{1.4321}{t}
Det nærliggande moglege er avgrensa. Det ligg spent ut mellom det busette på den eine sida og det forbodne og dei opne regionane lenger ute (1.431, 1.4311) på den andre.
\end{formproposition}

\begin{formproposition}{1.5}{d}
Ved sida av den materielle delinga går ei epistemisk.\kjelder{16,17}

Ho gjeld ein agent \(A\), eit omsyn \(k\), oppgjevne prøvevilkår og tidspunkt \(t\). Tidshorisonten høyrer til prøvevilkåra.

Posisjonar i \(\mathbb M\) der kriteriet langs \(X_k\) er avgjort under desse vilkåra, høyrer til det kjende \(K_{A,k,t}\); dei som oppfyller det, til \(K^+_{A,k,t}\).

Tenkte posisjonar i \(\mathbb M\) der kriteriet ikkje er avgjort, høyrer til det tenkte \(C_{A,k,t}\).

Resten er det utenkjelege langs omsynet:
\[
U_{A,k,t}=\mathbb{M}\setminus(K_{A,k,t}\cup C_{A,k,t}).
\]

Mengdene er parvis disjunkte og dekkjer det oppgjevne formrommet. Partisjonen er \(\Pi_{A,k,t}=(K_{A,k,t},C_{A,k,t},U_{A,k,t})\).

Der agent, vilkår og tidspunkt står faste, skriv eg \(K_k\), \(C_k\), \(U_k\) og \(\Pi_k\). Den samla kunnskapsstanden over dei utskilde omsyna er familien \(\Pi=\{\Pi_k\}\). Omsyn som enno ikkje er skilde ut (1.322), er ikkje eigne aksar i denne familien; dei høyrer til det akselause utanfor kartet.

Når omsynet òg står fast, nyttar eg dei korte nemningane \(K\), \(C\) og \(U\). Det akselause er ikkje ei fjerde delmengd i den aksevise partisjonen.
\end{formproposition}

\begin{formproposition}{1.501}{d}
\(C_k\) kallar eg tankerommet og \(K_k\) kunnskapsrommet langs \(X_k\). Rørsla i og mellom dei er den epistemiske sida av formgjevinga.
\end{formproposition}

\begin{formproposition}{1.51}{t}
Dei to delingane fell ikkje saman.

Ein forboden region kan vere tenkt: materielt utelukka, men tenkt.

Ein open region kan vere utenkjeleg: materielt tillaten, men utanfor det nokon har tenkt.
\end{formproposition}

\begin{formproposition}{1.511}{t}
Begge delingane trengst difor i ei skildring av formverda. Den materielle delinga seier kva vilkåra tillèt, den epistemiske kva som er kjent og tenkt.
\end{formproposition}

\begin{formproposition}{1.512}{t}
Grensene mellom K, C og U er epistemiske (1.2453). Dei flyttar seg når noko vert tenkt eller avgjort (1.2231).
\end{formproposition}

\begin{formproposition}{1.52}{t}
Det som er prøvd i ein realisert konfigurasjon, gjev kunnskap under prøvevilkåra. Slutningar frå det prøvde kan òg avgjere noko som enno ikkje er realisert.

\(K_k\) rommar både suksessar og fiaskoar. Ei avgjerd flyttar posisjonen til \(K_k\) berre for kriteriet og prøvevilkåra ho gjeld, ikkje for alle moglege verknader av forma. Same posisjon i \(\mathbb{M}\) kan difor liggje i \(K_k\) langs eitt omsyn og i \(C_j\) eller \(U_j\) langs eit anna (7.1).
\end{formproposition}

\begin{formproposition}{1.53}{t}
Ein kan kjenne at noko manglar utan å kjenne det som manglar. Innanfrå kan U nemnast som komplement. I det eit bestemt alternativ vert tenkt, forlèt det U.
\end{formproposition}

\begin{formproposition}{1.531}{t}
Frå ein vidare kunnskapsstand kan ein peike ut kva som låg i U for ein annan. Det same alternativet høyrer då til ulike delar av dei to kunnskapsstandane.
\end{formproposition}

\begin{formproposition}{1.5311}{d}
Det felles kjende for ei samling agentar langs \(X_k\) er posisjonane som ligg i \(K_{A,k,t}\) for kvar av dei, under prøvevilkår kvar av dei kan etterprøve. Ei avgjerd éin agent ikkje kan dele, utvidar hans eige \(K\), ikkje det felles.
\end{formproposition}

\begin{formproposition}{1.54}{t}
Den epistemiske rørsla kan ta utgangspunkt i det kjende \(K\) eller det tenkte \(C\), og utvide eller avgjere det eine ved hjelp av det andre.

\(C\to C\): Eit tenkt alternativ vert delt eller vidareført til nye alternativ.\newline{}\(C\to K\): Eit tenkt alternativ vert avgjort under det oppgjevne kriteriet.\newline{}\(K\to C\): Det kjende opnar eit spørsmål som enno ikkje er avgjort.\newline{}\(K\to K\): Ein kjend påstand vert prøvd eller presisert med finare oppløysing.

Dette er retningar i kunnskapsarbeidet, ikkje fire materielle operasjonar. Berre ei relevant prøve kan gjere eit tenkt alternativ kjent (1.52).
\end{formproposition}

\begin{formproposition}{1.541}{t}
Tanken veks innanfrå: frå det tenkte (C\(→\)C) eller frå det kjende (K\(→\)C). Det som kjem inn frå U, kjem ikkje ved at nokon ser over grensa (1.53).
\end{formproposition}

\begin{formproposition}{1.55}{d}
Den delen av \(C_k\) som ligg i det nærliggande moglege, er fronten langs \(X_k\): tilgjengeleg og tenkt, men ikkje kjend. Der kan det tenkte prøvast (1.52) og verte kjent.

Fronten er relativ til agenten, omsynet og vilkåra i den epistemiske delinga (1.5).
\end{formproposition}

\begin{formproposition}{1.6}{d}
Konfigurasjonar lukka under sum, differanse og produkt, og under transformasjonar \(θ\), med delrelasjonen \(≤\), utgjer ein formalgebra.\kjelder{11,13,18}
\end{formproposition}

\begin{formproposition}{1.61}{d}
I formalgebraen utfører formgrammatikken diskrete omskrivingar. Ein regel kan nyttast når venstresida finst som delform i den noverande forma, også når ingen på førehand namngav henne som eit objekt.

Ein formgrammatikk er \(SG=(S,R,s_0)\), der \(S\subseteq\mathbb{M}\) er formposisjonar, \(R=\{a_r\to b_r\}\) er reglar, og \(s_0\in S\) er startforma. Ei innleiring \(\theta\) utløyser eit steg når \(\theta(a_r)\leq s\):

\[
s\longrightarrow(s-\theta(a_r))+\theta(b_r).
\]

Eit steg tel berre når resultatet ligg i \(S\). Språket \(L(SG)=\{s\in S:s_0\to^{*}s\}\) er dei formene som kan genererast ved ei endeleg rekkje steg. Ein regelbruk svarar på geometrien slik han er no, ikkje på rekkjefølgja forma tidlegare vart bygd i.
\end{formproposition}

\begin{formproposition}{1.611}{t}
Formalgebraen ser gjennom det valde nivået, der objekta er atom. Der delane er linjer, flater og volum, har algebraen ingen atom, og innleiringa finn delformer som inga objektdeling gav på førehand.

Nivåbundne valensar uttømmer difor ikkje kva ei anna oppdeling kan gjere synleg (1.221, 1.223).
\end{formproposition}

\begin{formproposition}{1.62}{d}
Med \(SG\) meiner eg heretter ein historisk grammatikk: dei reglane som er kjende og i bruk. Språket hans er den delen av det oppgjevne formrommet som kan nåast frå startforma.

Nye reglar kan utvide språket eller krevje eit vidare rom. At det moglege er gjeve (1.222), garanterer ingen endeleg eller bereknbar grammatikk for alle moglege former.

Éin regelbruk er eitt steg. Det nærliggande moglege (1.432) ligg eitt steg frå det busette.
\end{formproposition}

\begin{formproposition}{1.63}{t}
Ein grammatikkoperasjon kan gjere ei moglegheit tenkt, frå ei tenkt eller kjend form. Han treng ikkje avgjere om forma vil verke.

Realisering er ei prøve på det tenkte. Ho vert ein \(C\to K\)-operasjon for eit oppgjeve kriterium når resultatet avgjer dette spørsmålet under prøvevilkåra.
\end{formproposition}

\begin{formproposition}{2}{a}
Under dei same vilkåra er ikkje alle moglege former like sannsynlege.
\end{formproposition}

\begin{formproposition}{2.01}{t}
Forbodet stengjer nokre alternativ. Fordelinga over dei tillatne krev i tillegg ein oppgjeven genereringsprosess eller ei avgrensa rekkje forsøk. Også eit tilfeldig utfall kan ha ei forklaring.
\end{formproposition}

\begin{formproposition}{2.011}{o}
Dei busette formene er ujamt fordelte innanfor det som var materielt tillate då dei vart til (1.4). Liknande former kan oppstå i tradisjonar utan dokumentert overføring. Likskapen åleine viser ikkje om han skuldast felles vilkår, overføring eller måten formrommet er valt på.
\end{formproposition}

\begin{formproposition}{2.02}{t}
Ei skeiv fordeling er noko som skal forklarast, ikkje sjølv ei forklaring. Forklaringa må vise kva som gjorde skilnad mellom alternativa (1.13).
\end{formproposition}

\begin{formproposition}{2.021}{t}
Seleksjonen endrar fordelinga gjennom det som skjer i kvart einskilt forløp. At ei form anten er realisert eller ikkje (1.12), hindrar ikkje at dei same prosessane flyttar, endrar og bryt henne.
\end{formproposition}

\begin{formproposition}{2.1}{d}
Eit seleksjonstrykk endrar sannsynsfordelinga over det som ikkje er forbode, relativt til ei oppgjeven grunnfordeling.\kjelder{3}

Trykket favoriserer og svekkjer alternativ utan å stengje dei. I landskapsspråket (3.1) er retninga mot favoriserte alternativ eit fall i spenninga.
\end{formproposition}

\begin{formproposition}{2.101}{t}
Forma er eitt realisert utfall av det trykka favoriserte, det agentane gjorde, og det vilkåra tillét. Dei andre alternativa treng ikkje vere utelukka då ho vart til.\kjelder{15,77,78}
\end{formproposition}

\begin{formproposition}{2.11}{t}
Eit trykk som ikkje tynner ut noko, er ikkje noko trykk. Samlar fordelinga seg i éin region, er utfallet bunde til regionen; først når éin posisjon står att, er heile utfallet determinert.
\end{formproposition}

\begin{formproposition}{2.111}{t}
Mellom nulltrykket og determinismen vert opne alternativ vekta ulikt utan at eitt utfall er tvinga fram. Her kan seleksjon og val skiljast frå både lik vekting og nødvendigheit.
\end{formproposition}

\begin{formproposition}{2.112}{a}
Design arbeider mellom desse grensene. Ei form kan likevel vere heilt bunden langs eitt omsyn og fri langs eit anna.
\end{formproposition}

\begin{formproposition}{2.12}{a}
Fleire seleksjonstrykk verkar samstundes. Korleis dei verkar saman, må skiljast frå at dei kan namngjevast kvar for seg.
\end{formproposition}

\begin{formproposition}{2.121}{t}
Eitt namngjeve trykk kan samle fleire omsyn og skjule avvegingane mellom dei. At modellen har éin verdi, viser ikkje at verda har eitt omsyn.
\end{formproposition}

\begin{formproposition}{2.122}{t}
Eit trykk som er likt over alle observerte tilfelle, er latent i samanlikninga. Det kan avgrense alle utan å forklare skilnaden mellom dei. Variasjon i vilkåra kan vise verknaden.
\end{formproposition}

\begin{formproposition}{2.13}{a}
Omsyn kjem i konflikt når dei ikkje kan møtast fullt ut i same tilgjengelege posisjon.\kjelder{120} Nokre konfliktar er reelle; andre kjem av eit for trongt repertoar av steg. Formgjevinga må skilje dei.
\end{formproposition}

\begin{formproposition}{2.131}{t}
Ei realisert form under ein slik konflikt er eit kompromiss: minst eitt relevant omsyn er ikkje fullt møtt.\kjelder{115} Utan konflikt kan fleire omsyn oppfyllast samstundes.
\end{formproposition}

\begin{formproposition}{2.1311}{t}
Utan eit oppgjeve sett av trykk er «suboptimal» tomt.
\end{formproposition}

\begin{formproposition}{2.14}{d}
Den delte moglegheita for operasjonar som avgrensar ein klasse (1.3), er klassen sin affordansestruktur.\kjelder{12,55,56}
\end{formproposition}

\begin{formproposition}{2.141}{t}
Affordansestrukturen er konstant innanfor klassen, og ein konstant forklarer ikkje åleine ein skilnad. Det klassen er til for, seier kvar formene samlar seg og kor trongt, ikkje kvifor dei er ulike.
\end{formproposition}

\begin{formproposition}{2.15}{t}
Ei god forklaring viser kva som skilde alternativa, og føreseier skilnader ho ikkje vart tilpassa til. Talet på namngjevne trykk tel mindre enn forklaringsarbeidet dei gjer.
\end{formproposition}

\begin{formproposition}{2.2}{d}
Materialaffordansen er operasjonane materialet tillèt og hindrar.\kjelder{12}

Utover forbodet i 1.42 er han ikkje ein regel, men ein gradient. Han bind sannsynsfordelinga utan å tvinge geometrien (2.1).

Signaturen hans i forma er latent (2.122) der alle tilfella har same materiale.
\end{formproposition}

\begin{formproposition}{2.201}{o}
Formhistoria er systematisk langsamare enn materialhistoria. Stål fanst i hundreår før stålrøyrstolen.\kjelder{54}

Først mangla operasjonen, og regionen var forboden (1.422). Så kom operasjonen, og regionen låg open, men ikkje tenkt (1.51).

Røyret låg i sykkelfabrikken i eit halvt hundreår. Nokon måtte sjå stolen i røyret.
\end{formproposition}

\begin{formproposition}{2.2011}{i}
Stolen låg i røyret som ei delform (1.611) ingen hadde skilt ut.
\end{formproposition}

\begin{formproposition}{2.202}{a}
Stilkategoriar som er fastsette uavhengig av formmåla dei skal føreseie, føreseier variasjon betre enn materiale innanfor ein klasse med liknande beregeometri.\kjelder{53}

Utval, aksar, stilklassifikasjon og grunnmodell må fastsetjast før prøva. Ei klassifisering av dei same formene ein vil predikere, er ikkje ein uavhengig forklaringsvariabel. Prøv difor modellane på former og tradisjonar som er haldne utanfor tilpassinga.
\end{formproposition}

\begin{formproposition}{2.2021}{t}
Materialet avgrensar kva som kan realiserast (2.2), men bestemmer ikkje åleine utfallet når fleire alternativ står opne. Det kviler på 2.12, ikkje på den uprøvde hypotesen i 2.202.
\end{formproposition}

\begin{formproposition}{2.21}{t}
Materialaffordansen er eitt trykk mellom mange. Ein samlevariabel kan predikere betre enn eit einskildtrykk ved å bere informasjon om fleire samtidige vilkår. Stilperioden kan vere ein slik variabel.
\end{formproposition}

\begin{formproposition}{2.211}{t}
Som samlevariabel er stilen eit symptom på trykka. Éin etikett kan samle fleire trykk utan å forklare korleis dei verka (2.02). Det krev ei eiga forklaring når stilen sjølv verkar på neste form (6.3).
\end{formproposition}

\begin{formproposition}{3}{t}
Seleksjonstrykka spenner ut eit landskap over formrommet.
\end{formproposition}

\begin{formproposition}{3.01}{a}
Lat \(\mu_0\) vere ei oppgjeven grunnvekt på dei tillatne formene. Dersom kvart trykk \(j\) multipliserer vekta i \(x\) med ein positiv faktor \(f_j(x)\), er den samla fordelinga \(\mu(x)\propto\mu_0(x)\prod_j f_j(x)\), føresett at normaliseringa er mogleg. Logaritmen av trykkfaktorane er då ein sum.

Dette er additivitetsvilkåret for landskapet. Der trykka samverkar ikkje-multiplikativt, må samspelsledda skiljast ut som eigne ledd i summen, eller faktorane definerast betinga av ei fastsett referanseordning.
\end{formproposition}

\begin{formproposition}{3.1}{d}
Tilpassingslandskapet \(Λ\) er den samla vektinga dei verksame trykka legg over formrommet under additivitetsvilkåret (3.01).\kjelder{3,8,58}

\(Λ\)(x, t) \(=\) \(∑_{j}\) \(λ_{j}\)(x, t), der \(λ_{j}\) \(=\) \(−\)ln \(f_{j}\). Verdien er spenninga i posisjonen. Berre skilnader i henne tel. Ved lik grunnvekt er lågare spenning meir sannsynleg.

Der gradientane finst, er \(−\)\(∇\)\(Λ\) \(=\) \(∑_{j}\) (\(−\)\(∇\)\(λ_{j}\)) den negative gradienten til den samla spenninga. Indeksen \(j\) omfattar både einskildtrykk og særskilt oppgjevne samspelsledd (3.01). Oppdelinga må følgje dei same faktorane og den same referanseordninga som summen. Der summen er ulik null, peikar han mot brattaste fall i \(Λ\) i det avstandsmålet som er valt på formrommet. Vert aksane skalerte om, endrar retninga seg; verdiane av \(Λ\) gjer det ikkje. Ved diskrete steg samanliknar me verdiane i naboane.
\end{formproposition}

\begin{formproposition}{3.11}{t}
Eit landskap lese ut av dei same formene det skal forklare, kan berre omnamne fordelinga. For å bere framover må vilkår eller overgangar fastsetjast uavhengig av det nye utfallet.
\end{formproposition}

\begin{formproposition}{3.12}{a}
Eit landskap kan ha fleire lokale minimum. Frå botnen av ein grunn dal kan alle dei tillatne stega føre opp, sjølv når det finst ein djupare dal på andre sida av ryggen.
\end{formproposition}

\begin{formproposition}{3.121}{d}
Eit lokalt minimum er ein posisjon utan betre tillatne nabosteg. Ein dal er ein region som ein bestemt navigasjonsmåte fører mot eit slikt minimum.

Dalen er ein attraktor for denne rørsla, ikkje for alle moglege rørsler.
\end{formproposition}

\begin{formproposition}{3.122}{t}
At eit landskap kan ha fleire lokale minimum (3.12), gjev ingen rett til å føresetje éi allmenn, fullkomen form.\kjelder{58} Eit optimum må alltid namngje klasse, tillatne steg, vilkår, tid og mål. Utan desse er «optimal» ikkje ein prøvbar påstand.
\end{formproposition}

\begin{formproposition}{3.2}{t}
Bratte dalveggar kan avgrense spreiinga; høge ryggar kan gjere dalen vanskeleg å forlate. Stabiliteten krev i tillegg ei rørsle som svarar på forstyrringar.

Når rørsla fører forstyrra former tilbake, er dei kanaliserte.\kjelder{8}
\end{formproposition}

\begin{formproposition}{3.21}{t}
Djupe, bratte kanalar fryser dimensjonane; slake rommar fluktuasjonen.
\end{formproposition}

\begin{formproposition}{3.3}{d}
Ein stilart er ei vedvarande samling av former som ligg nær kvarandre i eit oppgjeve empirisk formrom.\kjelder{9,53} Kva som er nært, kjem an på aksane og målet ein nyttar (1.321–1.322).

Stilen er ikkje sjølv eit objekt. Han er eit mønster i fordelinga av objekt. Stilperiodar kan overlappe og gli over i kvarandre; ein kategori som vert lagd på formene etterpå, er ikkje i seg sjølv ei forklaring.

Ein ser ikkje ein stil utan å sjå dei formene som ber likskapen.
\end{formproposition}

\begin{formproposition}{3.31}{t}
Kvar stol er bestemt, men svermen er eit mønster over mange stolar. Grensene må fastsetjast gjennom likskapen som vert undersøkt, ikkje gjennom kor bestemte stolane er.

Det som i statistikken er eit aggregat (2.211), kan i landskapet vise seg som ein sverm.
\end{formproposition}

\begin{formproposition}{3.4}{a}
Formgjeving arbeider under vilkår ho ikkje åleine rår over: materielle verkemåtar, andre agentar og behova formene skal møte.

Ei realisert form kan sjølv verte eit vilkår for nye former. Ho tek plass, bind materiale, går inn i nye konfigurasjonar og kan endre kva som vert føretrekt eller mogleg (1.412). Kor stor verknaden er, må undersøkast; ei realisering endrar ikkje nødvendigvis det målte landskapet.
\end{formproposition}

\begin{formproposition}{3.41}{t}
Uavhengig konvergens kan prøve hypotesen om ein felles, reell attraktor.\kjelder{117} Han styrkjer forklaringa når felles opphav og måten rommet er konstruert på, ikkje er nok til å forklare likskapen.
\end{formproposition}

\begin{formproposition}{3.411}{d}
Reell tyder: gjeven ved vilkåra i 3.4, naturlovene og dei behova formene skal møte, same kva tradisjon som navigerer.
\end{formproposition}

\begin{formproposition}{4}{a}
Landskapet ligg ikkje fast. Det vert endra av nye vilkår og av formene som vert realiserte i det.
\end{formproposition}

\begin{formproposition}{4.01}{t}
Formhistoria er difor dynamikk, ikkje berre fordeling. Når ei realisering endrar vilkår, tilgang eller seleksjonsvekt (3.4), endrar ho landskapet for seinare steg. Berre det at ei form er ny, er ikkje bevis på ei målbar endring i trykket.
\end{formproposition}

\begin{formproposition}{4.1}{a}
Seleksjonstrykka kan endre styrke og retning, også utan at nokon endrar mål med vilje.

Dalar kan verte djupare eller grunnare, flytte seg, dele seg og smelte saman. Ei tilpassing er difor relativ til dei vilkåra og det tidspunktet ho vert prøvd under.
\end{formproposition}

\begin{formproposition}{4.11}{o}
Ein ny teknologi flyttar grensa for det forbodne; ein ny marknad endrar kva landskapet belønner; eit nytt materiale utvidar operasjonsrommet.
\end{formproposition}

\begin{formproposition}{4.12}{t}
Gradvise trykkforskyvingar kan hope seg opp til ein dal forsvinn. Landskapet endrar seg jamt; staden for minimumet gjer eit sprang.

Ein slik overgang er ein bifurkasjon.\kjelder{111}
\end{formproposition}

\begin{formproposition}{4.2}{o}
I nokre formhistorier følgjer korte periodar med rask endring etter lange periodar med små synlege endringar. Eit slikt forløp viser ikkje åleine om skiftet kom av nye vilkår, nye reglar eller endra val.
\end{formproposition}

\begin{formproposition}{4.21}{t}
Når vilkåra endrar seg jamt og eit minimum forsvinn, er det brå formskiftet eit merke etter den jamne dynamikken (4.12). Andre sprang kan kome av brå vilkårsendringar eller val mellom dalar som framleis finst.
\end{formproposition}

\begin{formproposition}{4.3}{t}
Landskapet har minne.\kjelder{14}

Ei realisert form kan endre trykka for den neste (3.4), og kvar ny form tek utgangspunkt i ein busett posisjon (1.411). All formgjeving er difor omformgjeving: kvar ny form arvar eit landskap og ein posisjon.

Stiavhengigheit er eit grunntrekk, ingen anomali.
\end{formproposition}

\begin{formproposition}{4.4}{t}
Med kvar realisering endrar det busette seg, og det nærliggande moglege kan flytte seg.\kjelder{13,59,60,79}

Materialstraumar, nye operasjonar og læring kan opne vegar; tap av stoff, handlag og tilgang kan stengje dei. Det tilgjengelege veks berre der opninga er større enn tapet.
\end{formproposition}

\begin{formproposition}{4.5}{t}
Eitt trykk kan ha fleire dalbotnar. Fleire trykk kan møtast i éin. Av talet på dalar kan ein difor ikkje lese kor mange trykk som forma landskapet.\kjelder{59}

Kva dal ei form hamnar i, kjem òg an på kva vegar som stod opne då ho vart til (4.3).
\end{formproposition}

\begin{formproposition}{5}{a}
Formverda rommar agentar som svarar på landskapet.
\end{formproposition}

\begin{formproposition}{5.01}{t}
Landskapet flyttar inga form. Det verkar gjennom materialet, agentane og overgangane mellom konfigurasjonar.
\end{formproposition}

\begin{formproposition}{5.011}{t}
Ein formtype som går att under endra vilkår, kan verte halden gjennom regulering, eller vere stabil utan henne. Å skilje dei krev å undersøkje kva som skjer når vilkåra vert forstyrra.
\end{formproposition}

\begin{formproposition}{5.012}{o}
Nokre formtypar går att gjennom skiftande vilkår. Stolar vert laga på nytt rundt kroppar med mykje av den same beregeometrien. Haldet må finnast i arbeidet som fører forma vidare.
\end{formproposition}

\begin{formproposition}{5.013}{t}
Å svare er å ta påverknad inn og gje påverknad ut. Analysen må difor skilje eit system frå omgjevnaden og finne vegane mellom dei. Grensa kan vere ein fysisk membran eller ei funksjonell avgrensing.
\end{formproposition}

\begin{formproposition}{5.02}{d}
Ei hinne er ei grense av tilstandar mellom eit indre \(I\) og eit ytre \(E\). I ein oppgjeven stokastisk modell er grensa eit Markov-teppe \(B\) når det indre \(I\) og det ytre \(E\) er betinga uavhengige gjeve \(B\).\kjelder{83}

Ei statistisk grense er ikkje nødvendigvis ein fysisk membran. Observasjon og handling går gjennom dei koplingane grensa tillèt. Sjølvhald krev i tillegg at systemet regulerer vilkåra for si eiga grense (5.021).
\end{formproposition}

\begin{formproposition}{5.021}{d}
Ei hinne er sjølvhalden i den grad regulering innanfor henne motverkar forstyrringar som elles ville løyse henne opp. Ei grense som berre den som analyserer, teiknar, er ei analysegrense.
\end{formproposition}

\begin{formproposition}{5.022}{d}
Eit sjølv er eit system avgrensa av ei sjølvhalden hinne. I ontologien i 1.24 er det eit mønster i feltet.
\end{formproposition}

\begin{formproposition}{5.0221}{o}
Hinner vert haldne i ulik grad: nokre av det indre, nokre av andre, dei fleste av begge.
\end{formproposition}

\begin{formproposition}{5.0222}{t}
Å vere eit sjølv er difor ein grad (5.021).
\end{formproposition}

\begin{formproposition}{5.023}{t}
Steinen held ei grense utan å regulere henne. Eit sjølv må motverke forstyrringar som elles ville løyse grensa opp. Begge følgjer naturlovene; skiljet er reguleringa.
\end{formproposition}

\begin{formproposition}{5.024}{a}
Den sjølvhaldne avgrensinga er òg ei avgrensing innanfrå: sjølvet er ei oppleving som ikkje når inn i andre mønster i feltet. Det ho har av resten, kjem gjennom hinna.

Utanfrå er sjølvet hinna og det ho held. Den regulerande og den opplevde avgrensinga høyrer til same organisering.
\end{formproposition}

\begin{formproposition}{5.025}{t}
Det sjølvet har av det ytre, er tilstandane i hinna. Det ser difor ikkje feltet, men ei koding av feltet på hinna.

Persepsjonen er ikkje eit vindauge, men eit grensesnitt.
\end{formproposition}

\begin{formproposition}{5.0251}{t}
Sjølvet har ikkje eit ukoda møte å samanlikne kodinga med. Det kan prøve henne gjennom det ho lèt systemet skilje, føreseie og halde.

Ei koding kan tene dette arbeidet utan å likne det ho kodar, og samstundes skjule noko systemet treng.
\end{formproposition}

\begin{formproposition}{5.026}{t}
Fleire sjølv i same felt har eit felles ytre, alt som ligg utanfor dei alle. Dei deler difor éi verd, sjølv om kvart har si eiga koding av henne.

Uavhengig konvergens (3.41) kan vere eit spor av den verda dei deler.
\end{formproposition}

\begin{formproposition}{5.027}{t}
Det eit sjølv vil, er innanfor hinna. Det verkar utover berre gjennom hinna, som handling. (Ønske rører ikkje fjellet; ei hand kan.)
\end{formproposition}

\begin{formproposition}{5.028}{t}
Kroppen til eit sjølv er korleis ei avgrensa oppleving tek seg ut utanfrå (5.024, 1.2351).

Biletet følgjer det det er bilete av, og lagar det ikkje.
\end{formproposition}

\begin{formproposition}{5.0281}{o}
Hos menneske kan endringar i hjernen følgje endringar i rapportert oppleving tett. Dette viser samvariasjon mellom eit ytre målbart mønster og ei opplevd side, ikkje at oppleving berre finst i hjernen eller at samvariasjonen i seg sjølv forklarer henne.
\end{formproposition}

\begin{formproposition}{5.1}{a}
I målretta formgjeving vert landskapet navigert av system med tilbakekopling: dei registrerer avvik og endrar verksemda som svar. Kor mykje eit bestemt forløp krev av slik agens, må visast.\kjelder{62}
\end{formproposition}

\begin{formproposition}{5.11}{d}
Ein agent er eit system som registrerer avvik frå ein målregion og har ein politikk som svarar for å minske avviket eller halde det innanfor toleransen.

Formelt er \(A=(m,\pi,G,\varepsilon)\). Modellen \(m\) knyter observasjonar, handlingar og venta følgjer saman; \(G\) er målregionen; \(\pi\) er politikken agenten tek i bruk; \(\varepsilon\) er toleransen. Ordninga treng ikkje vere medviten eller representert som eit bilete.

Lat \(e_k(t;\pi,d)\) vere avstanden frå målregionen langs omsynet \(X_k\), ved tid \(t\) under politikken \(\pi\) og forstyrringa \(d\). Avstanden og eininga må fastsetjast for kvart omsyn. \(\mathcal{P}_A\) er politikkområdet agenten faktisk har tilgang til. For ei realisert form \(x\) skriv eg \(e_k(x)\) for det største avviket over forstyrringane og tidsrommet i ei oppgjeven prøve.

Prediksjonsfeilen er ein annan storleik. Når observasjon og prediksjon er sannsynsfordelingar, kan han målast som \(D_k=D_{\mathrm{KL}}(p_{\mathrm{obs}}\Vert p_m)\).\kjelder{112} Låg \(D_k\) tyder at modellen treff det observerte; låg \(e_k\) tyder at utfallet møter målet. Eit godt estimat kan difor føreseie ein skade presist.

Agens krev denne svarstrukturen, ikkje medvit, nervesystem eller erklært intensjon. At svaret lukkast, er eit anna spørsmål.
\end{formproposition}

\begin{formproposition}{5.111}{t}
Kva agenten er laga av, avgjer ikkje om han er ein agent. Kvart system som oppfyller vilkåra, er det.
\end{formproposition}

\begin{formproposition}{5.1111}{i}
Bioelektriske mønster, muskelrørsler og språklege modellutdata kan alle inngå i reguleringssløyfer.\kjelder{22,24,85} Ingen av mekanismane er åleine bevis på agens. Berre der registrert mål-avvik faktisk endrar ein politikk, er definisjonen i 5.11 oppfylt.
\end{formproposition}

\begin{formproposition}{5.112}{t}
Variasjonell fri energi kan gje ei øvre grense for overraskinga i det observerte.\kjelder{19,112,114} Låg prediksjonsfeil er likevel ikkje låg skade: ein kan føreseie samanbrotet presist utan å kunne hindre det.
\end{formproposition}

\begin{formproposition}{5.113}{t}
Negativ tilbakekopling gjev ein agent når svaret er kopla til eit avvik frå det som skal haldast. Målet er regionen G, avviket er e, og justeringa er politikken som motverkar det.

Termostaten har eit band rundt ein temperatur. Brytaren koplar målinga til varme som reduserer avviket.\kjelder{4,5,7} Han treng ikkje føreseie rommet langt fram for å halde det no.
\end{formproposition}

\begin{formproposition}{5.114}{d}
Haldshorisonten \(τ_{k}\)(A) er kor lenge ein tilgjengeleg politikk kan halde \(e_{k}\) \(≤\) \(ε_{k}\) under oppgjevne forstyrringar. Prøva må òg fastsetje starttilstanden og eventuell innstillingstid.

Prediksjonshorisonten \(\tau^{\mathrm{pred}}_k(A)\) er kor lenge modellen føreseier følgjene langs \(X_k\) med oppgjeven presisjon. Ein agent kan ha lang prediksjonshorisont og kort haldshorisont, eller halde noko ved lokal tilbakekopling som han ikkje føreseier langt fram.
\end{formproposition}

\begin{formproposition}{5.1141}{a}
Ein realisert modell har endeleg mange parametrar.
\end{formproposition}

\begin{formproposition}{5.1142}{t}
Horisonten må prøvast under forstyrringane oppgåva krev. Endeleg parametrisering avgjer ikkje åleine kvar han sluttar; lokal regulering kan vare medan andre sider av verda er ukjende.
\end{formproposition}

\begin{formproposition}{5.115}{t}
Eit sjølv, sett utanfrå, er ein agent som held si eiga hinne. Reguleringa svarar på avvik som truar vilkåra for hennar eiga vidare verksemd. Dette er sjølvhald (5.021, 5.11).
\end{formproposition}

\begin{formproposition}{5.116}{t}
Ikkje kvar agent er eit sjølv.

Termostaten held ein temperatur, men hinna hans vert halden av nokon annan: han vert bygd og halden ved like utanfrå.

Han er ein agent, ikkje eit sjølv: hinna vert halden heilt utanfrå (5.0222).
\end{formproposition}

\begin{formproposition}{5.117}{t}
Strukturen legg ikkje oppleving til. Det strukturen er bilete av (1.2351), er opplevinga.
\end{formproposition}

\begin{formproposition}{5.118}{t}
Den som veit og har tenkt i 1.5, er ein agent. Kunnskapsstanden er alltid kunnskapsstanden til nokon som navigerer.
\end{formproposition}

\begin{formproposition}{5.12}{o}
Agens er lagdelt: strukturen er stabil, oppløysinga varierer.
\end{formproposition}

\begin{formproposition}{5.121}{i}
Termostaten: eitt mål, éin dimensjon, éin brytar.
\end{formproposition}

\begin{formproposition}{5.122}{i}
Treet: fleire gradientar, ingen sentral prosessor. Kvar grein navigerer lokalt.
\end{formproposition}

\begin{formproposition}{5.123}{i}
Kolonien: millionar av agentar. Reiret trer fram i feltet mellom agentane.
\end{formproposition}

\begin{formproposition}{5.124}{i}
Pattedyret: prediksjon, minne\kjelder{101}, planlegging.\kjelder{64,65} Representasjonen strekkjer seg over tid.
\end{formproposition}

\begin{formproposition}{5.125}{i}
Den som formgjev for andre: representasjonen strekkjer seg forbi behova til agenten sjølv, til dei som skal leve med forma.
\end{formproposition}

\begin{formproposition}{5.13}{a}
Agenten føretrekkjer låg spenning langs dei omsyna måla gjeld. Han svarar på avvik slik dei når han gjennom hinna, og vel etter dei følgjene modellen hans ventar.

Å gå ned i estimert mål-avvik kan då føre mot målet. Det fører ned i det verkelege landskapet berre i den grad estimatet er rett og det som er valt, kan gjennomførast.\kjelder{19,64}
\end{formproposition}

\begin{formproposition}{5.131}{t}
Målet er ein region i eit gjeve rom, men agenten kan endre kva region han held som mål. Det moglege vert ikkje laga av ønsket; retninga til arbeidet kan verte det.
\end{formproposition}

\begin{formproposition}{5.132}{t}
Lokal navigasjon kan gje global form utan at nokon ser heilskapen.\kjelder{19} Men lokale fall kan stengje vegen til eit betre mål. Kompetanse kan krevje å gå opp ei stund for å kome ned ein annan stad.
\end{formproposition}

\begin{formproposition}{5.14}{d}
Spenninga i 3.1 ligg over formrommet. Den indre spenninga \(σ_{k}\)(A) er avstanden langs \(X_{k}\) mellom det agenten møter og målregionen hans, slik han verkar inne i agenten. Ho set politikken i rørsle.\kjelder{118}
\end{formproposition}

\begin{formproposition}{5.141}{t}
Det som set politikken i rørsle, kjem frå minst tre stader: verda mot målet, sansinga mot forventinga, og målet langs éin akse mot målet langs ein annan. Den første er mål-avviket e, som \(σ_{k}\) i 5.14 gjeld. Den andre er prediksjonsfeilen D. Den tredje er ein konflikt mellom omsyn.

For å minske e ved eit fast mål kan agenten handle gjennom \(π\) og endre situasjonen. Ved høg D kan han prøve føreseiinga mot nye observasjonar og rette modellen m (5.112). Når omsyn strir, kan han finne eit steg som møter dei saman, eller avgjere kva som må vike (2.13, 5.32).
\end{formproposition}

\begin{formproposition}{5.142}{t}
Ein bjelke, ei celle og eit menneske kan alle vere i spenning. Storleikane vert ikkje identiske av namnet. I ein agent må det visast korleis avviket vert registrert og set ein politikk i gang.
\end{formproposition}

\begin{formproposition}{5.143}{t}
Utan registrert spenning kan målet vere nådd, gløymt eller ikkje lenger sansa. Ro er difor ikkje åleine eit bevis på at det som skal haldast, er trygt.
\end{formproposition}

\begin{formproposition}{5.144}{a}
For eit sjølv er den indre målspenninga kjend innanfrå som uro når det registrerte avviket står over toleransen; når det fell tilbake under, som ro.

Uro treng ikkje vere medviten eller uttrykkje den verkelege faren presist. Ho er innsida av svaret, ikkje ein ufeilbarleg rapport om verda.
\end{formproposition}

\begin{formproposition}{5.1441}{o}
Svolt, smerte og frykt følgjer truslar mot det ein organisme må halde. Metting, lindring og tryggleik følgjer at trusselen fell bort.
\end{formproposition}

\begin{formproposition}{5.15}{d}
Evna til å nå same målregion frå ulike utgangspunkt og på ulike vegar er kompetanse.
\end{formproposition}

\begin{formproposition}{5.151}{t}
Kompetansen vert synleg når den vanlege vegen vert sperra: kva andre vegar kan agenten finne til same mål?
\end{formproposition}

\begin{formproposition}{5.152}{d}
Kompetansen kan målast som stegsparing mot eit oppgjeve mål G, samanlikna med ei bestemt blind leiting:\kjelder{86}

\(ν\)(A; G) \(=\) \(\mathrm{log}_{10}\)(E \(N_{\mathrm{blind}}\)(G) / E \(N_{A}\)(G)).

Forventingane gjeld same fordeling av utgangspunkt og forstyrringar, same stegeining og same avbrotsgrense, og må vere endelege og positive. Positiv \(ν\) er gevinst, negativ \(ν\) tap. Målet seier kva prøva kostar i steg, ikkje kva G er verdt.
\end{formproposition}

\begin{formproposition}{5.153}{t}
Kvar måling av kompetanse er gjord av ein observatør, i eit rom observatøren har valt (1.321). Målinga seier difor ikkje berre noko om agenten, men òg om observatøren.
\end{formproposition}

\begin{formproposition}{5.2}{d}
At avvik set ein agent i arbeid for å måle, endre eller halde noko, er omhug \(ω\)(A).\kjelder{23,68}

\(ω_{k}\)(A) er styrken på svaret langs \(X_{k}\) under oppgjevne vilkår. For å samanlikne styrken må ein fastsetje både eininga for avviket og kva som tel som arbeid.

Omhug er ein profil over omsyn, ikkje eitt tal for kor mykje agenten bryr seg om alt.
\end{formproposition}

\begin{formproposition}{5.201}{t}
Omhug viser seg i prioritering av arbeid. Kunnskap åleine endrar ikkje det kjende; omhug åleine gjer ikkje eit feil estimat rett. Innsats, handlag og utfall må difor prøvast kvar for seg.
\end{formproposition}

\begin{formproposition}{5.21}{d}
Kapasiteten \(\kappa(A)\) er dei omsyna agenten har ein tilgjengeleg politikk for å halde innanfor toleransen, over eit oppgjeve tidsrom og under oppgjevne forstyrringar.\kjelder{21,86}

For ei prøve med tidsrom \(T\) og forstyrringar \(\mathcal D\):

\[
\begin{aligned}
\kappa(A;T,\mathcal D)=\{X_k:\ &\exists\pi\in\mathcal P_A\;\forall d\in\mathcal D,\\
&\forall t\in T:\ e_k(t;\pi,d)\leq\varepsilon_k\}.
\end{aligned}
\]

Dette er kapasitet langs kvar akse for seg. Den samla kapasiteten er ei mengd av aksemengder:

\[
\begin{aligned}
\kappa_{\mathrm{samla}}(A;T,\mathcal D)=\{S\subseteq\mathcal X:\ &\exists\pi\in\mathcal P_A\;\forall X_k\in S,\\
&\forall d\in\mathcal D\;\forall t\in T:\\
&e_k(t;\pi,d)\leq\varepsilon_k\}.
\end{aligned}
\]

Her er \(\mathcal X\) dei oppgjevne omsynsaksane. Den same politikken må halde alle aksane i \(S\). Difor er \(X_k\in\kappa\) det same som \(\{X_k\}\in\kappa_{\mathrm{samla}}\), medan \(S\subseteq\kappa\) ikkje åleine gjev \(S\in\kappa_{\mathrm{samla}}\).

Starttilstanden og ei eventuell innstillingstid må vere fastsette før prøva. Å kjenne ein politikk ein ikkje kan ta i bruk, er ikkje å ha han i kapasiteten.
\end{formproposition}

\begin{formproposition}{5.211}{t}
Kapasitetsprøva gjev eit hald agenten kan setje i verk. Prediksjonsprøva gjev eit samband han har grunn til å vente. Eit vellukka resultat i den eine prøva kan ikkje telje som resultat i begge.
\end{formproposition}

\begin{formproposition}{5.22}{d}
Det som kan setje agenten i spenning, er den kognitive lyskjegla \(Γ\)(A): aksane der \(ω_{k}\) \(>\) 0, kvar med så lang utstrekning fram og tilbake i tid som måla hans rekk langs han.\kjelder{21,85,90,173}

Lyskjegla er ikkje grensa for kva agenten sansar eller når, men for kva han har som mål.

\(I_{Γ}\)(A) \(=\) \{k: \(ω_{k}\)(A) \(>\) 0\} er akseprojeksjonen av lyskjegla. \(τ^{Γ}_{k}\)(A) seier kor langt målet langs \(X_{k}\) rekk fram i tid, rekna frå same starttid som haldsprøva.
\end{formproposition}

\begin{formproposition}{5.221}{i}
Lyskjegler spenner frå cellulær millisekundrespons til tiårslang konsensus.
\end{formproposition}

\begin{formproposition}{5.222}{t}
Det som ligg utanfor lyskjegla, set ikkje målretta arbeid i gang. Det kan likevel vere kjent, og det kan verke på agenten anten han kjenner det eller ikkje.
\end{formproposition}

\begin{formproposition}{5.223}{t}
Lyskjegla avgjer kva skilnader i formrommet som har noko å seie for agenten.

Innanfor henne møter han ikkje former (5.025), men operasjonane konfigurasjonen tillèt han, affordansar.\kjelder{12,55}

Den delen av dei som er felles for heile klassen, er affordansestrukturen i 2.14; det agenten navigerer etter, er resten.
\end{formproposition}

\begin{formproposition}{5.2231}{t}
Ei vidare lyskjegle gjer latente affordansar verksame; ei trongare gjer verksame latente.\kjelder{55,70}
\end{formproposition}

\begin{formproposition}{5.23}{t}
Lyskjegla og kapasiteten treng ikkje falle saman. Ein agent kan bry seg om det han ikkje kan halde, og kunne halde det han ikkje bryr seg om.
\end{formproposition}

\begin{formproposition}{5.231}{d}
\(H_0(A)\) omfattar aksane som ligg både i lyskjegla og i kapasiteten kvar for seg, under dei same prøvevilkåra.

\[
H_0(A)=\{X_k\in\kappa(A):\omega_k(A)>0\}.
\]

To aksar kan liggje i \(H_0(A)\) kvar for seg utan å kunne haldast samstundes. Eit samla hald under éin politikk \(\pi\in\mathcal P_A\) er

\[
\begin{aligned}
H(A;\pi)=\{X_k\in H_0(A):\ &e_k(t;\pi,d)\leq\varepsilon_k\\
&\text{for alle }d\in\mathcal D,\ t\in T\}.
\end{aligned}
\]

For politikken \(\pi\) agenten faktisk tek i bruk, skriv eg \(H(A)=H(A;\pi)\). Dette er det robuste haldet under dei oppgjevne prøvevilkåra, og \(H(A;\pi)\in\kappa_{\mathrm{samla}}(A;T,\mathcal D)\). Om politikken held aksane, må prøvast.
\end{formproposition}

\begin{formproposition}{5.232}{d}
Ein akse kan falle utanfor \(H_0(A)\) på tre måtar.

Utan omhug og utan kapasitet: blindskap.

Med kapasitet, utan omhug: likegyldigheit.

Med omhug, utan kapasitet: avmakt.

Blindskap er her ei operativ nemning: ho seier at agenten verken svarar eller kan halde, ikkje at han nødvendigvis manglar all kunnskap om aksen.
\end{formproposition}

\begin{formproposition}{5.2321}{t}
Vilkåra verkar vidare under blindskap, likegyldigheit og avmakt. Utfallet kan verte verna eller skada gjennom andre omsyn og gjennom substratet; det uhaldne manglar eit eige hald.
\end{formproposition}

\begin{formproposition}{5.2322}{t}
Ein akse kan ikkje haldast lenger enn haldshorisonten \(τ_{k}\)(A) rekk. Lyskjegla er ikkje bunden av denne horisonten. Når \(τ^{Γ}_{k}\)(A) \(>\) \(τ_{k}\)(A), er tidsrommet mellom dei to horisontane avmakt i tid, under dei oppgjevne prøvevilkåra. Det gjeld òg når agenten kan føreseie kva som vil skje.
\end{formproposition}

\begin{formproposition}{5.233}{t}
Ein målregion må skilje mellom tilstandar. Ein region som femner alle alternativ likt, gjev inga retning. Å halde den tilstanden agenten alt er i, kan derimot krevje vedvarande arbeid.
\end{formproposition}

\begin{formproposition}{5.234}{t}
I gapet mellom omhug og samla hald bur strevet. Sjølv der kvart omsyn kan haldast for seg, kan oppfinninga stå att. Ei ny form eller eit nytt steg kan gjere det mogleg å halde omsyna saman (2.13, 5.231, 5.32).

Først når éin tilgjengeleg politikk held alle desse omsyna samstundes under prøvevilkåra, er det arbeidet med å halde dei vidare som står att (4.1, 6.511).
\end{formproposition}

\begin{formproposition}{5.24}{t}
Omhugen avgjer korleis lyskjegla er fylt: tett langs nokre aksar, tynn langs andre.

Spesialist og generalist er to profilar av same slag. Spesialisten har få høge komponentar, generalisten mange lågare. Samanlikninga gjeld den relative profilen (6.61), ikkje den samla styrken.
\end{formproposition}

\begin{formproposition}{5.241}{t}
Ein agent med trong omhug kan ha stor kapasitet: mange aksar innanfor rekkjevidd, få som set noko i gang.

Likegyldigheita (5.232) er då ikkje per agent, men per akse.
\end{formproposition}

\begin{formproposition}{5.25}{t}
Omhugen kan ikkje fyllast ut i ettertid for ei form som alt er navigert. Komponentane måtte vore der då forma vart til.
\end{formproposition}

\begin{formproposition}{5.26}{t}
Kor langt eit sjølv rekk, er kor langt lyskjegla rekk.

Hinna seier kvar sjølvet møter verda; lyskjegla seier kva i verda som har noko å seie for sjølvet.

Rekkjevidda krympar, veks og flyttar seg når måla gjer det; hinna kan stå.
\end{formproposition}

\begin{formproposition}{5.261}{t}
For eit sjølv er omhug, sett innanfrå (5.024, 5.144), at noko har noko å seie. Sett utanfrå er omhug (5.2) kor mykje arbeid avvik set i gang i politikken.
\end{formproposition}

\begin{formproposition}{5.2611}{t}
Det som utanfrå er registrert mål-avvik over toleransen langs ein akse i lyskjegla, er innanfrå uro (5.144). At ein annan måler den faktiske faren, er noko anna enn at sjølvet registrerer han.
\end{formproposition}

\begin{formproposition}{5.2612}{t}
Omhug er difor ikkje ei kjensle i tillegg til politikken. Det same svaret som vert undersøkt utanfrå, er kjent innanfrå som at noko har noko å seie.
\end{formproposition}

\begin{formproposition}{5.262}{t}
Eit omsyn som ikkje vert registrert som målrelevant, utløyser ikkje målretta uro. Men ein kan kjenne at ein ikkje bryr seg. Kunnskap om ei grense er ikkje det same som å flytte henne.
\end{formproposition}

\begin{formproposition}{5.263}{t}
Agenten kan oppdage nokre av grensene sine utan å vite at han har funne alle. Eiga gransking, forstyrringar og andre sine svar kan gjere ulike utelatingar synlege.\kjelder{70}
\end{formproposition}

\begin{formproposition}{5.27}{d}
Kompetanse i eit rom er intelligens: evna til å finne vegar til målregionen, òg vegar som fører bort frå han ei stund og nærare seinare.\kjelder{23,86,91}

Omhug er innhaldet i intelligensen. Omhugen gjev intelligensen noko å gjelde; intelligensen gjev omhugen vegar.
\end{formproposition}

\begin{formproposition}{5.271}{t}
Meir omhug endrar kva agenten prioriterer; meir handlag kan endre kva han maktar. Dei to verknadene kan prøvast ved å endre den eine medan den andre vert halden fast.
\end{formproposition}

\begin{formproposition}{5.272}{t}
Intelligensen gjeld berre innanfor lyskjegla. Eit problem som ikkje kan setje agenten i spenning, er ikkje hans problem.
\end{formproposition}

\begin{formproposition}{5.28}{d}
At avviket hos ein annan agent set agenten i arbeid for å minske det for den andre, er omsorg. Omsorg er den relasjonelle utvidinga av målrommet, aksane målregionen \(G\) er fastsett langs (5.11): lyskjegla får aksar som høyrer til ein annan.\kjelder{23,68,90}

Med omsorg er vurderinga til \(A\) (5.7)

\[
\begin{aligned}
\Phi_A={}&(1-\gamma_A)\sum_k\omega_k(A)\hat\Lambda_{A,k}\\
&+\sum_{B\ne A}\sum_j c_{A\to B,j}\hat\Lambda_{A|B,j}.
\end{aligned}
\]

\(\hat\Lambda_{A|B,j}\) er biletet \(A\) har av \(e_{B,j}/\varepsilon_{B,j}\): avviket til \(B\) frå sin eigen målregion langs \(X_j\), målt i toleransen hans. Målregionen og toleransen er den andre sine. Biletet er \(A\) sitt.

\(A\) har omsorg for \(B\) langs \(X_j\) når \(c_{A\to B,j}>0\). Omsorgsvektene summerer seg til omsorgsdelen \(\gamma_A\leq1\); omhugsprofilen \(\omega(A)\) fordeler resten på agenten sine eigne omsyn. Omsorg er omhug (5.2) langs aksane til andre.
\end{formproposition}

\begin{formproposition}{5.281}{t}
Omsorg er ikkje eit nytt slag svar. Ho er omhug langs ein akse der målregionen og toleransen høyrer til ein annan, og vert prøvd på same måte: ved arbeidet avviket set i gang, ikkje ved det agenten seier.

Ein slik akse ligg i kapasiteten når \(A\) har ein tilgjengeleg politikk som held \(e_{B,j}\leq\varepsilon_{B,j}\) under prøvevilkåra (5.21). Politikken til \(B\) høyrer då til prøvevilkåra.
\end{formproposition}

\begin{formproposition}{5.2811}{d}
Vekta til ein akse i vurderinga til \(A\) er \((1-\gamma_A)\omega_k(A)\) for eit eige omsyn og \(c_{A\to B,j}\) for eit omsyn hos ein annan. Der 5.22, 5.231 og 5.232 krev \(\omega_k(A)>0\), gjeld kravet denne vekta.
\end{formproposition}

\begin{formproposition}{5.282}{d}
Ein agent kan gjere tre ting med avviket hos ein annan: føreseie det, bruke det og ta det inn som mål.

Å føreseie er å ha avviket i modellen \(m\). Å bruke er å la det tene eit mål som ikkje er den andre sitt (6.021). Å ta det inn som mål er å ha det i sitt eige målrom. Berre det siste er omsorg.
\end{formproposition}

\begin{formproposition}{5.2821}{t}
Avhengnad er ikkje omsorg. Når eit eige omsyn hos \(A\) samvarierer med avviket hos \(B\), fordi \(B\) held noko \(A\) treng, kan arbeidet til \(A\) verne \(B\) utan omsorgsvekt (6.675). Det er bruk, og vernet rekk så langt som avhengigheita.

Omsorg er det som står att der avviket hos den andre ikkje lenger tener agenten sjølv (6.701).
\end{formproposition}

\begin{formproposition}{5.283}{d}
Langs eigne aksar kan agenten ta feil av verda. Langs aksane til ein annan kan han òg ta feil av målet.

Omsorga verkar når tre vilkår held: vekta er positiv, biletet følgjer det faktiske avviket hos den andre (6.71), og aksen ligg i kapasiteten.

Utan vekt og utan kapasitet: blindskap. Kapasitet utan vekt: likegyldigheit. Vekt utan kapasitet: avmakt (5.232).

Vekt og kapasitet med eit bilete som ikkje følgjer avviket: feilretta omsorg. Kjem feilen av at agenten tek feil av verda, er han den same som langs eigne aksar (6.674). Kjem han av at agenten tek feil av målet, er det formynderi: arbeid for eit mål den andre ikkje har.
\end{formproposition}

\begin{formproposition}{5.284}{t}
Omsorga flyttar grensa for intelligensen (5.272): det som set den andre i spenning, vert eit problem for agenten. Ho gjev intelligensen meir å gjelde. Vegane må agenten framleis finne.
\end{formproposition}

\begin{formproposition}{5.3}{t}
Det agenten navigerer i, er det grammatikken (1.62) kan generere.
\end{formproposition}

\begin{formproposition}{5.301}{t}
Grammatikken avgrensar det genererbare, og innanfor det vektar landskapet (2.1) vegane.
\end{formproposition}

\begin{formproposition}{5.31}{t}
Av det grammatikken tilbyr, realiserer agenten helst det som minskar spenninga (5.13) langs aksane i lyskjegla hans.
\end{formproposition}

\begin{formproposition}{5.32}{t}
Dei fire operatorane i 1.54 kan alle endre kva politikkar agenten har tilgang til. Ein ny tenkt framgangsmåte kan opne eit steg; ei prøve kan vise om steget verkar.

Kapasiteten vert ikkje større berre fordi noko nytt er tenkt eller kjent. Han vert større når eit nytt hald kan gjennomførast.
\end{formproposition}

\begin{formproposition}{5.4}{d}
Substratet er dei materielle vilkåra formgjevinga arbeider gjennom og mot. Det gjev motstand og resonans og legg ned veto mot nokre overgangar.\kjelder{26,81,92,93}

Som «agent av nullte orden» markerer det grensa der berre den materielle verknaden står att, utan modell, politikk og omhug.
\end{formproposition}

\begin{formproposition}{5.401}{o}
I modellar av kompetente materiale kan me skilje startkonfigurasjonen frå dei reglane som endrar henne.\kjelder{92} I visse simuleringar gjer endring av regelparametrar navigasjonen meir robust under støy enn endring av startforma. Det er ein modellspesifikk observasjon, ikkje ei lov om alle materiale.
\end{formproposition}

\begin{formproposition}{5.4011}{t}
Strukturparametrane og funksjonsparametrane svarar til startforma og reglane i ein formgrammatikk (1.61).
\end{formproposition}

\begin{formproposition}{5.402}{t}
Substratet bidreg på to måtar: hardt, ved å leggje ned veto og avgrense kvar landskapet gjeld (1.42), og mjukt, som eitt trykk (2.2) i landskapet.
\end{formproposition}

\begin{formproposition}{5.403}{t}
Substratet kan bere mekanisk spenning utan omhug. Rørsla følgjer verkemåtane i materialet. Ho tel som navigasjon hos ein agent berre der eit registrert avvik set ein politikk i gang.
\end{formproposition}

\begin{formproposition}{5.5}{o}
Evna til å navigere hopar seg opp som strukturert informasjon.\kjelder{64}
\end{formproposition}

\begin{formproposition}{5.51}{d}
Læring er ei endring i agenten sin modell eller politikk som følgje av erfaring, og som gjer ein skilnad i ny navigasjon. Læring kan utvide, snevre inn eller omfordele kapasiteten. Det som vert bevart i eit felles arkiv, treng ikkje verte bevart i kvar agent.
\end{formproposition}

\begin{formproposition}{5.511}{t}
Læring kan utvide det tilgjengelege: eit nytt hald opnar nye steg. Kvar realisering gjer ein posisjon busett; ho vert ny kunnskap når prøva avgjer noko agenten ikkje visste (1.63).
\end{formproposition}

\begin{formproposition}{5.52}{t}
Læring kan utvide \(\kappa(A)\), men garanterer ikkje omhug. Agenten kan få fleire tilgjengelege måtar å halde noko på utan å prioritere dei (5.21–5.23).
\end{formproposition}

\begin{formproposition}{5.521}{t}
Ei lære som gjer eitt omsyn til det einaste, kanaliserer arbeidet etter dette omsynet. Dei andre trykka verkar vidare, medan læra sluttar å krevje svar på dei.\kjelder{15,78}

Dette er ei innsnevring av omhugen, ikkje berre ei matematisk samanfatting av fleire omsyn i éin verdi.
\end{formproposition}

\begin{formproposition}{5.5211}{t}
Der det utelatne krev eit svar politikken ikkje gjev, vert manglande hald ei kjelde til svikt. Kva følgjer utelatinga får, vert føresagt gjennom dei verksame vilkåra (7.152).
\end{formproposition}

\begin{formproposition}{5.5212}{t}
Der læra rår, vert handlaget som svara på dei andre trykka, ikkje lenger øvt.

Det lèt seg ikkje gjenskape frå læra, for éin dimensjon ber ikkje informasjonen til dei andre.

Då krympar det tilgjengelege, og handlaget er tapt i agenten (5.51), sjølv om det står i arkivet.
\end{formproposition}

\begin{formproposition}{5.53}{t}
Med agentar får landskapet ein ny kanal for endring.

Det kjende går inn i modellane som styrer kva agentane realiserer (5.11, 5.31), og det som endrar sannsynsfordelinga over det som ikkje er forbode, er trykk (2.1).

Det kjende verkar difor sjølv som trykk (1.52, 3.4) på det som kjem etter.
\end{formproposition}

\begin{formproposition}{5.6}{d}
Liv kallar eg den sjølvhaldande prosessen som opprettheld vilkåra for si eiga vidare målretta regulering, og dermed for si eiga lyskjegle.\kjelder{85,89} Kjegla kan utvidast, men utviding er ikkje eit vilkår for liv. Eit system med mål som vert halde heilt av andre, oppfyller ikkje denne definisjonen ved målet åleine (5.115–5.116).
\end{formproposition}

\begin{formproposition}{5.61}{t}
Det levande motverkar at vilkåra for eiga vidare regulering bryt saman. Kor sterkt sjølvhaldet er, må prøvast under forstyrring. Breidda på lyskjegla åleine er ikkje eit mål på om systemet lever.
\end{formproposition}

\begin{formproposition}{5.62}{t}
Det levande utøver sjølvhald (5.115) og er difor eit sjølv sett utanfrå, også når det er avhengig av andre for å halde seg.
\end{formproposition}

\begin{formproposition}{5.63}{d}
Død er eit irreversibelt samanbrot i evna til å halde den same levande avgrensinga ved like. Når denne prosessen ikkje kan takast opp att, fell også lyskjegla til dette sjølvet bort.
\end{formproposition}

\begin{formproposition}{5.631}{t}
Fråvær av synleg målretta rørsle i eit tidsrom viser ikkje at lyskjegla er tapt. Ein må skilje ein kvilande prosess frå ein som har stansa irreversibelt. Døden ligg på grensa for det sjølvhaldet som definerer liv (5.6).
\end{formproposition}

\begin{formproposition}{5.632}{t}
Døden avsluttar ei bestemt avgrensing, ikkje feltet. At feltet held fram, inneber ikkje at den avgrensa opplevinga held fram.
\end{formproposition}

\begin{formproposition}{5.6321}{t}
Kva det avgrensa var for seg sjølv, kan ikkje kjennast utanfrå (1.235). Kva som held fram av det, kan difor ikkje avgjerast her.
\end{formproposition}

\begin{formproposition}{5.64}{t}
Grensa mellom levande og ikkje-levande er ikkje ei grense mellom to slag substans, for same felt er på begge sider (1.24).\kjelder{89}

Skilnaden er kor mykje av hinna det indre held sjølv (5.0222). Han er ikkje i slag, men i grad.
\end{formproposition}

\begin{formproposition}{5.641}{a}
System som held seg ved å svare, kan modellerast ved prediksjon og regulering.\kjelder{89,112} Liv krev at denne reguleringa held vilkåra for eiga vidare verksemd.
\end{formproposition}

\begin{formproposition}{5.7}{a}
Kjernelikninga seier korleis ein agent vel mellom tilgjengelege grammatikksteg.\kjelder{11,17,21} Lat \(q_A(x,x')\) vere sannsynet for å føreslå eit tillate steg frå \(x\) til \(x'\). Ein lokal overgangskjerne kan då skrivast som

\[
\begin{aligned}
T_A(x,x')={}&q_A(x,x')\min\bigl\{1,\\
& r_A(x,x')e^{-\beta[\Phi_A(x')-\Phi_A(x)]}\bigr\},\\
&\quad x'\ne x .
\end{aligned}
\]

\[
\begin{aligned}
\Phi_A(x)&=\sum_{k\in I_\Gamma(A)}
  \omega_k(A)\hat\Lambda_{A,k}(x),\\
\hat\Lambda_{A,k}(x)&=
  \hat e_{A,k}(x)/\varepsilon_{A,k}.
\end{aligned}
\]

Her er \(\hat e_{A,k}\) agenten sitt estimerte mål-avvik langs omsyn \(k\), \(\varepsilon_{A,k}>0\) toleransen og \(\omega_k(A)\) dei normaliserte omhugvektene (6.61). \(\beta\geq0\) styrer kor sterkt estimerte skilnader vektar vala. Har agenten omsorg, kjem ledda i 5.28 til.

Ved direkte lokal utforsking kan ein setje \(r_A=1\). Det gjev ikkje åleine ei kjend stasjonær fordeling. Når ein i staden vil ha ei reversibel utforsking med grunnfordeling \(\mu_{SG}\), krevst Hastings-korreksjonen

\[
r_A(x,x')=
\frac{\mu_{SG}(x')q_A(x',x)}
{\mu_{SG}(x)q_A(x,x')}.
\]

Korreksjonen gjeld for forslag med positivt sannsyn begge vegar; utan retursteg vert aksepten null (6.64). Resten av overgangssannsynet vert verande i \(x\). Materielle steg som ikkje kan gjerast om att, krev difor ei eiga skildring (6.641).

Grammatikken avgrensar kva \(q_A\) kan føreslå. \(\kappa(A)\) seier noko anna: kva agenten kan halde under forstyrring. Ei form kan endrast utan at målet vert nådd. Modellen føreseier ei fordeling over steg, ikkje den verkelege følgja av dei.

Lat forslaga vere symmetriske med kovarians \(\Sigma_q\), inne i eit glatt tilgjengeleg rom med lokalt flat grunnfordeling. Når stega er små mot \(1/(\beta|\nabla\Phi_A|)\), er gjennomsnittsdrifta per steg \(-\frac{\beta}{2}\Sigma_q\nabla\Phi_A\) til første orden. Ved større steg eller brattare fall flatar drifta ut og veks ikkje lenger med \(\beta\). Det er ei tilnærming, ikkje ei generell rørslelov. For endeleg \(\beta\) kan stega òg gå oppover.
\end{formproposition}

\begin{formproposition}{5.71}{d}
Kjernen i 5.7 skil mellom forslaga \(q_A\), den estimerte vurderinga \(\Phi_A\), seleksjonsstyrken \(\beta\) og dei materielle vilkåra. Desse ledda må ikkje tilpassast til det same utfallet som skal brukast til å stadfeste modellen.

Med \(r_A=1\) får me fordelinga i 6.6 når forslaga har detaljert balanse med grunnfordelinga og dei andre vilkåra i 6.64 held. Hastings-korreksjonen gjev same fordeling for fleire reversible forslag. Vel agenten likt mellom dei \(N(x)\) stega som kan gjerast om att frå \(x\), er \(\mu_{SG}(x)\propto N(x)\): former med mange opne steg veg tyngst i grunnfordelinga. Ingen av desse konstruksjonane gjer ei irreversibel materialhandling reversibel.
\end{formproposition}

\begin{formproposition}{5.72}{t}
Reduksjon.\kjelder{11,17} Fire grensetilfelle viser kva komponentane gjer:

(i) Utan seleksjon står dei faktisk tilgjengelege grammatikkstega og grunnfordelinga deira att.

(ii) Eit landskap som berre skil det kjende vellukka K\(+\) frå resten, gjev ei binær vekting av kjent og uavgjort. Det skildrar ikkje åleine korleis nye tankar vert til.

(iii) Med fast grammatikk, gjeve parameterrom og agentar som kan styre alle aksane i det, får me optimering i eit fast rom.

(iv) Med rette estimat, full kapasitet og omhug langs alle omsyn fell uvissa om vurdering og gjennomføring bort. Konflikt mellom mål kan framleis stå att.
\end{formproposition}

\begin{formproposition}{5.721}{t}
Formgrammatikken gjev det generative laget, tanke-kunnskapsskiljet det epistemiske laget, og agentanalysen laget for mål og gjennomføring. Dei tre laga er generative grammatikkar, kunnskapsutvikling og agens på fleire skalaer. Samanbindinga forklarer korleis kvart lag avgrensar og opnar dei andre.
\end{formproposition}

\begin{formproposition}{5.722}{a}
Historisk formgjeving bind dei tre laga saman. Ei avgrensa undersøking kan halde eitt lag fast; eit fundament for faget må òg forklare overgangane mellom dei.
\end{formproposition}

\begin{formproposition}{5.7221}{t}
Den allvitande agenten er eit grenseomgrep. Ei konkret forklaring krev at kunnskap og kapasitet vert prøvde, ikkje føresette som uavgrensa.
\end{formproposition}

\begin{formproposition}{5.723}{t}
Dei tre laga stiller ulike forklaringskrav: korleis forma kan genererast, korleis det tenkte vert kjent, og korleis vurdering vert verksam i handling. Å løyse eitt av krava er ikkje å ha løyst dei andre.
\end{formproposition}

\begin{formproposition}{5.7231}{a}
Fullstendigheitshypotesen er at dei tre laga saman ber dei tre ledda ei forklaring av formgjeving treng: generering, vekting og målretta gjennomføring. Grammatikken ber genereringa. Vektinga ligg i landskapet \(\Lambda\), som agentane og vilkåra ber saman (6.613). Gjennomføringa ligg i kapasiteten til agentane. Kunnskapsstanden avgrensar kva som kan føreslåast, og kor rett det vert vurdert.

Hypotesen fell dersom eit nødvendig ledd har eigne, etterprøvbare verknader som ikkje kan fangast av nokon av dei tre utan at verkemåten går tapt. Å gje feilen eit nytt namn i eitt av ledda er ikkje å forklare han. Kvar kjend årsak får plassen sin i eitt ledd før prøva. Å flytte henne etter eit motdøme er eit tap (7.173).
\end{formproposition}

\begin{formproposition}{5.724}{t}
Formrommet \(\mathbb M\) er grunnen i modellen. Kunnskapsstanden \(\Pi=(\Pi_{A,k,t})\) er familien av epistemiske delingar relativt til agent, kriterium og tidspunkt, under oppgjevne prøvevilkår (1.5).

Femtuppelet \((\mathbb M,\Pi,SG,\Lambda,\{A_i\})\) held frå kvarandre grunnen \(\mathbb M\), dei tre laga \(SG\), \(\Pi\) og \(\{A_i\}\), og landskapet \(\Lambda\) som vektar alternativa, utan å gjere dei til isolerte ting.
\end{formproposition}

\begin{formproposition}{6}{t}
Forma oppstår i samspelet mellom grammatikk, kunnskap og agens på fleire skalaer.
\end{formproposition}

\begin{formproposition}{6.001}{t}
Kvar navigert form har ein som navigerer og eit substrat som gjer navigasjonen mogleg og avgrensar henne. Fleire agentar kan føre kvar sin politikk inn i same form.

Det ein agent bryr seg om, men ikkje kan halde, krev hjelp, meir kapasitet eller endra vilkår.
\end{formproposition}

\begin{formproposition}{6.002}{t}
Feltet mellom agentane er same felt som sjølva er mønster i. Hinnene avgrensar tilgangen innanfrå, ikkje eit anna stoff. Det som ligg utanfor kvar hinne, kan vere felles vilkår for fleire.
\end{formproposition}

\begin{formproposition}{6.003}{t}
Trykka har både materielle og agentborne verkemåtar. At ei avgrensing kjem frå ein politikk, gjer henne ikkje verksam utan stoff og energi. Skiljet gjeld korleis ho vert halden og kan endrast.
\end{formproposition}

\begin{formproposition}{6.01}{d}
Ein agent som ved å navigere realiserer konfigurasjonar i formrommet, er ein formgjevar.

Cella, bia, treet, handverkaren og ingeniøren er formgjevarar. Det som varierer, er materialet agensen er bygd av, ikkje strukturen.
\end{formproposition}

\begin{formproposition}{6.02}{d}
Ein designar er ein formgjevar som endrar vilkåra ein annan agent navigerer under (3.4), og har omsorg for han (5.28): eit avvik hos den andre set arbeid i gang for å minske avviket for den andre.

Kriteriet gjeld bestemte omsyn. Det krev eit arbeid som kan rettast av det som held eller sviktar hos den andre (6.712), ikkje at arbeidet alltid lukkast.
\end{formproposition}

\begin{formproposition}{6.021}{t}
Verknad åleine er ikkje nok. Den som berre bruker den andre sine avvik til å styre han mot eigne mål, formgjev vilkåra hans, men møter ikkje designkriteriet langs dette omsynet.
\end{formproposition}

\begin{formproposition}{6.0211}{t}
Den som endrar vilkåra, treng ikkje bere følgjene eller gje den andre same moglegheit til å endre dei tilbake. Denne asymmetrien er situasjonen det faglege ansvarsprinsippet (7.26) skal svare på.
\end{formproposition}

\begin{formproposition}{6.022}{t}
Designkriteriet er operasjonelt og kan prøvast også utanfor yrkesrolla. Ei celle som regulerer verksemda si etter mål-avviket hos ein nabo, kan langs dette omsynet oppfylle kriteriet. Det tilskriv ikkje cella menneskeleg moralsk ansvar.

Eit slikt svar krev meir enn påverknad; den andre sitt avvik må endre politikken. Ingen er designar langs alle omsyn fordi kriteriet er møtt langs eitt.
\end{formproposition}

\begin{formproposition}{6.023}{t}
Å vere designar er eit spørsmål om grad: kor mange aksar og kor mange andre lyskjegla femner.

Omgrepet vert meir alvorleg, ikkje mindre. Det gjeld langt fleire system enn profesjonen har rekna med.

Profesjonen sjølv er designar berre langs dei aksane der kriteriet er oppfylt. Kva aksar det er, må han vise.
\end{formproposition}

\begin{formproposition}{6.0231}{t}
Kor mykje ein designar formgjev for, er ikkje kor langt omsorga hennar når i rom, men kor mange ulike målregionar til andre lyskjegla hennar femner.
\end{formproposition}

\begin{formproposition}{6.0232}{t}
Ingen er nøydd til å vere designar. Definisjonen pålegg ingenting. Han fastset kva namnet krev.

Den som kallar arbeidet sitt design langs eit omsyn, seier noko som kan prøvast (6.022) og kan vere usant.
\end{formproposition}

\begin{formproposition}{6.024}{t}
Omsorga tek situasjonen til den andre inn i det ein sjølv prøver å halde. Biletet av måla hans kan vere feil; arbeidet må kunne rettast av eit svar ein ikkje alt har bestemt innhaldet i.
\end{formproposition}

\begin{formproposition}{6.025}{t}
Å formgje for ein annan er å gjere avviket hos den andre til eit kriterium i eige arbeid. Avviket og biletet ein har av det, må haldast frå kvarandre.

Observasjon og handling kan bere omsyn over ei grense utan å oppheve grensa mellom den som formgjev, og den som lever med forma.
\end{formproposition}

\begin{formproposition}{6.026}{t}
Ein agent kan leggje spenninga si over på ein annan ved å omsetje henne til eit teikn den andre handlar på. (Bestillinga er eit slikt teikn; forma er svaret.)

Overføringa krev ein felles protokoll. Utan han er teiknet støy.
\end{formproposition}

\begin{formproposition}{6.027}{t}
Formgjeving endrar vilkåra andre møter, også langs omsyn som låg utanfor lyskjegla til den som formgav (5.222). Asymmetrien i 6.0211 gjeld difor òg følgjer ingen hadde som mål å skape (7.01, 7.21).
\end{formproposition}

\begin{formproposition}{6.1}{a}
Former i samansett praksis er polyagentiske kompromiss: fleire agentar og materielle prosessar verkar inn på kva som vert til.\kjelder{66,67}

Korleis kontrollen er fordelt, er eit spørsmål om verknad, ikkje berre om kor mange som var til stades. Ein agent kan dominere utan å rå over alle følgjene.
\end{formproposition}

\begin{formproposition}{6.11}{a}
Agens er organisert på fleire skalaer: delar kan regulere noko, heilskapen noko anna.\kjelder{22,67,85,87} Kvar skala må ha den svarstrukturen han skal forklare.
\end{formproposition}

\begin{formproposition}{6.111}{a}
Heilskapen styrer delane ved å endre vilkåra dei svarar på. Slik kan lokale mål tene eit mål ingen del representerer.
\end{formproposition}

\begin{formproposition}{6.1111}{t}
Eitt forløp kan skildrast både som lokale kausale overgangar og som regulering på eit høgare nivå. Skalaen vert forklarande når han gjev eigne føreseiingar om kva heilskapen held under forstyrring. Å kalle eit mønster målretta utan å finne ein svarstruktur er ikkje å forklare det (5.11).
\end{formproposition}

\begin{formproposition}{6.112}{t}
Kvart nivå navigerer innanfor si eiga lyskjegle (5.272).\kjelder{67} Delane på lågare nivå kan difor fremje mål langs aksar som ligg i lyskjegla til heilskapen (6.111), men ikkje i deira eiga.
\end{formproposition}

\begin{formproposition}{6.113}{t}
Forklaringa nedanfrå viser korleis delane verkar; forklaringa ovanfrå viser kva felles hald verksemda inngår i. Ei god fleirskalaforklaring må binde dei saman i dei same overgangane.\kjelder{72}
\end{formproposition}

\begin{formproposition}{6.114}{t}
Når fleire agentar både føreslår og prøver nye steg, kan rørsla mellom tanke og kunnskap vere distribuert på tvers av skalaer (1.54, 1.63). Det inneber ikkje at kvar materialendring er læring, eller at alle delane har same kunnskapsstand.
\end{formproposition}

\begin{formproposition}{6.12}{a}
Overgangar i individualitet gjer spørsmålet om skala konkret: når kan mange agentar telje som éin?\kjelder{87}

Koplinga må gje delane eit felles avvik dei kan regulere. Lik rørsle kan oppstå også når delane svarar kvar for seg.
\end{formproposition}

\begin{formproposition}{6.121}{d}
Interaksjonsstrukturen er djup når spenninga vert delt: avviket i den eine vert eit avvik i den andre, utan merke om kven det høyrer til.\kjelder{173}
\end{formproposition}

\begin{formproposition}{6.1211}{t}
Når eit avvik vert eit felles reguleringssignal, kan arbeidet hos éin del minske avviket hos ein annan. Delen oppfyller då designkriteriet langs dette omsynet dersom den andre sitt mål vert halde som kriterium (6.02).
\end{formproposition}

\begin{formproposition}{6.1212}{t}
Når avviket er delt, kan eit lokalt svar halde noko som fleire delar er avhengige av.
\end{formproposition}

\begin{formproposition}{6.1213}{t}
Der skaden på éin del vert ført tilbake som eit avvik i den felles reguleringa, kan heilskapen svare på han. Utan slik tilbakeføring kan éin del vinne på tapet til ein annan.
\end{formproposition}

\begin{formproposition}{6.1214}{t}
Delt spenning er grensetilfellet av gjensidig omsorg (6.741): omsorg utan adresse. I omsorga veit agenten kven avviket høyrer til. I den djupe koplinga set avviket arbeid i gang same kvar det oppstod.
\end{formproposition}

\begin{formproposition}{6.122}{a}
Koplinga må bere det felles haldet og samstundes bevare dei skilnadene arbeidet krev. Det optimale ligg i dette forholdet, ikkje i sterkast mogleg kopling.
\end{formproposition}

\begin{formproposition}{6.1221}{t}
Når ei kopling som bar eit felles sjølvhald, brest, kan delane halde fram utan at heilskapen gjer det. Det som døyr på det høgare nivået, treng ikkje døy på det lågare.
\end{formproposition}

\begin{formproposition}{6.123}{d}
Når ei samling agentar oppfyller 5.11 som heilskap, er ho éin agent.

Det er eit skalasprang når heilskapen held det ingen del held: ein akse utanfor lyskjeglene til alle delane, eller ein lengre horisont enn nokon av delane.
\end{formproposition}

\begin{formproposition}{6.1231}{t}
Delt spenning (6.121) gjev samlinga ei felles registrering av avvik. Det 5.11 krev i tillegg av heilskapen, er ein målregion og ein toleranse som gjeld heilskapen, og ein politikk som faktisk held det felles avviket innanfor toleransen (6.132).

Delt spenning er difor nødvendig for at mange skal telje som éin langs eit omsyn (6.12), men ikkje tilstrekkeleg. Lik vurdering (6.74) viser at samlinga vel som éin. At ho held som éin, må prøvast under forstyrring (5.21, 6.1111).
\end{formproposition}

\begin{formproposition}{6.13}{t}
Agenten vel og verkar; vilkåra opnar, hindrar og prøver. Kontroll og seleksjon kan vere verksame både før, under og etter realiseringa.\kjelder{66}
\end{formproposition}

\begin{formproposition}{6.131}{o}
Utforsking krev lause koplingar; utnytting krev tette.
\end{formproposition}

\begin{formproposition}{6.132}{t}
Kompetente delar kan rette eit avvik før det vert svikt i heilskapen.\kjelder{116} Delt spenning må då møte kapasitet til å svare.
\end{formproposition}

\begin{formproposition}{6.1321}{t}
Ein plan treng difor ikkje vere fullstendig for å verke. Der materialet sperrar vegen (5.402), rettar handverkaren det teikninga ikkje såg.
\end{formproposition}

\begin{formproposition}{6.1322}{t}
Kompetansen i delane gjev heilskapen tid til å prøve det nye utan å døy av det.
\end{formproposition}

\begin{formproposition}{6.14}{t}
Designaren verkar aldri åleine (6.1).\kjelder{66,67,93} Ho møter andre agentar som har eigne mål, og materielle prosessar som ikkje treng å ha mål. Kompetansen ligg ikkje i å late som om alle vil det same, men i å finne former som kan halde fleire omsyn under dei faktiske vilkåra.
\end{formproposition}

\begin{formproposition}{6.2}{o}
Når nokre delar vert bundne saman og andre ikkje, oppstår tilbakekoplingssløyfer som går gjennom heilskapen.\kjelder{66}
\end{formproposition}

\begin{formproposition}{6.201}{t}
Ei global tilbakekopling bind dei lokale sløyfene saman når følgjene av svara deira kjem tilbake som eit felles reguleringssignal (6.1213, 6.123).
\end{formproposition}

\begin{formproposition}{6.21}{o}
Hinner ligg i hinner.

Cella held si, vevet si om cellene, organismen si om vevet.
\end{formproposition}

\begin{formproposition}{6.211}{t}
Når hinnene mellom dei lågare sjølva ber avvik over, og ei ny sjølvhalden hinne vert halden om alle, oppstår eit sjølv av høgare orden. Dette er ei ny avgrensing i feltet, ikkje ei summering av opplevingar.

Det som må forklarast, er korleis den felles reguleringa vert halden, og kva tilgang den nye avgrensinga gjev.
\end{formproposition}

\begin{formproposition}{6.3}{t}
Utanom agentane har ein stilart inga kausal kraft (2.211).\kjelder{9,15,53}

I den første verkstaden bøygde snikkaren stolryggane etter malen. I den neste skar ein annan snikkar malen etter bogen ho kjende att i stolane.

Likskapen mellom stolane var eit spor etter arbeidet (2.211, 3.3). Då snikkaren lét bogen styre snittet, fekk sporet verknad i det neste arbeidet (5.53). Ein kan slik føre ein stil vidare utan å kjenne trykka som gav han form.
\end{formproposition}

\begin{formproposition}{6.31}{t}
For brukaren kan stilarten verke som eit grensesnitt, ikkje berre som eit spor etter seleksjonstrykk.\kjelder{80} Brukaren møter attkjennbare former, medan trykka som forma dei, ofte ikkje er synlege.
\end{formproposition}

\begin{formproposition}{6.311}{t}
Designaren som ignorerer grensesnittet, ignorerer agenten ho formgjev for.\kjelder{80}
\end{formproposition}

\begin{formproposition}{6.312}{t}
Stilen er eit grensesnitt i same meining som persepsjonen er det (5.025, 5.0251): ei koding som viser det brukaren treng for å navigere, ikkje trykka bak. Stilen er kodinga ein kultur deler, persepsjonen kodinga til ein kropp.
\end{formproposition}

\begin{formproposition}{6.32}{t}
Å navigere med kunnskap om trykka bak stilen, ikkje berre om tyngdepunktet hans, er å representere fleire aksar.
\end{formproposition}

\begin{formproposition}{6.4}{t}
Ei realisert form utvidar det kjende når møtet med henne avgjer eit nytt spørsmål, og vert eit utgangspunkt for neste steg.\kjelder{14,79}
\end{formproposition}

\begin{formproposition}{6.5}{t}
Inga form har krav på å vere den siste. Nye vilkår kan endre kva som verkar, og nye omsyn kan vise at det som verka, hadde ein kostnad.\kjelder{60,61}

Former og grammatikkar kan vare lenge. Det læra skal bevare gjennom endringane, er evna til å undersøkje korleis reglar, vilkår og agentar ber forma.
\end{formproposition}

\begin{formproposition}{6.51}{t}
Ein omhug som veks raskare enn kapasiteten, opnar nye oppgåver før dei gamle er haldne. Han kan drive arbeidet vidare sjølv når eit einskilt verk er ferdig.
\end{formproposition}

\begin{formproposition}{6.511}{t}
Ein agent kan ha som mål å utvide måla sine. Då kan ei måloppnåing opne ei ny oppgåve.

Også eit uendra mål kan krevje arbeid utan siste steg: vedlikehaldet sluttar ikkje fordi det lukkast. Å opne og å halde er ulike former for vidare verksemd.
\end{formproposition}

\begin{formproposition}{6.5111}{t}
Å opne nye omsyn krev å finne nye prøver og vegar. Å halde noko kjent krev å møte det som forstyrrar det. Kor mykje arbeid kvar av dei krev, avheng av vilkåra.
\end{formproposition}

\begin{formproposition}{6.6}{a}
Den allmenne forma bind saman det genererbare, det vurderte og det gjennomførte. Fordelingsmodellen er:\kjelder{86,89}

P(x \(|\) \{\(A_{i}\)\}, SG, V) \(=\) \(Z^{-1}\) \(μ_{\mathrm{SG}}\)(x) exp(\(−\)\(β\) \(Φ\)(x)), for x \(∈\) D.

\(Φ\)(x) \(=\) \(∑_{i}\) \(w_{i}\) \(\hat{Λ}_{i}\)(x).

Her er D ei ikkje-tom mengd former som kan nåast frå startforma med tillatne steg utanfor V. \(μ_{\mathrm{SG}}\) er ei oppgjeven grunnfordeling med positiv vekt på D. Estimata og vektene står faste.

For eit endeleg D med endelege verdiar gjev maksimalentropiprinsippet fordelinga direkte som valmodell: ho maksimerer den relative entropien \(-\sum_{x\in D}P(x)\log(P(x)/\mu_{SG}(x))\) under ei øvre grense for det forventa vekta avviket \(\Phi\). For ei grense strengt over minimumet er \(\beta\geq0\) multiplikatoren til grensa; når ho ikkje bind, er \(\beta=0\). Grensa ved minimumet vert nådd som \(\beta\to\infty\) (6.612).\kjelder{113}
\end{formproposition}

\begin{formproposition}{6.61}{d}
\(\hat{Λ}_{i}\)(x) \(=\) \(∑_{k∈I_{Γ}(A_{i})}\) \(ω_{k}\)(\(A_{i}\)) \(\hat{Λ}_{\mathrm{ik}}\)(x), der \(\hat\Lambda_{ik}(x)=\hat e_{ik}(x)/\varepsilon_{ik}\) er agenten sitt estimerte mål-avvik langs omsyn k, skalert med toleransen \(\varepsilon_{ik}>0\) (5.11). Estimatet kan avhenge av heile konfigurasjonen; verdien 1 svarar til toleransegrensa.

\(w_i\geq0\) er agenten sin relative innverknad, med \(\sum_iw_i=1\). I denne valmodellen er \(\omega_k(A_i)\geq0\) den relative omhugsprofilen, normalisert med \(\sum_{k\in I_\Gamma(A_i)}\omega_k(A_i)=1\) når lyskjegla er ikkje-tom. Med omsorg (5.28) går dei eigne omsyna inn med faktoren \(1-\gamma_i\), og omsorgsvektene summerer til \(\gamma_i\); \(\hat\Lambda_i\) er då vurderinga \(\Phi_{A_i}\). Ved tom lyskjegle set eg \(\hat\Lambda_i=0\) og \(w_i=0\). \(\beta\geq0\) er seleksjonsstyrken. Vektene, dei toleranseskalerte estimata og \(\beta\) er dimensjonslause og må fastsetjast før modellen vert prøvd; kva vala åleine kan fastsetje, står i 6.6732.

Produktet \(\theta_{ik}=\beta w_i\omega_k(A_i)\) er den absolutte svarstyrken i 5.2, målt i log-odds per toleranseeining.

Lyskjegla fastset vurderinga, kapasiteten kva vurderinga kan gjennomføre. \(μ_{\mathrm{SG}}\) kjem frå både grammatikken og den oppgjevne stegprosessen.
\end{formproposition}

\begin{formproposition}{6.611}{t}
Fordelinga gjeld vala under oppgjevne vilkår. Dei faktiske følgjene krev ei eiga skildring av møtet mellom forma og verda.

Materielle vilkår verkar òg før realiseringa ved å endre kva steg som er tilgjengelege, og kor sannsynlege dei er.
\end{formproposition}

\begin{formproposition}{6.612}{t}
Ved \(β\) \(=\) 0 står grunnfordelinga att. Når \(β\) veks, aukar vekta til posisjonane med låg \(Φ\) relativt til dei med høg. Når \(Φ\) når eit globalt minimum i D, samlar massen seg i grensa \(β\) \(→\) \(∞\) om dei globale minima, fordelte etter grunnvekta dersom fleire minimum står likt.

Sterk seleksjon er ikkje nødvendigvis eitt utfall. Ho kan heller ikkje gjere eit feil estimat rett.
\end{formproposition}

\begin{formproposition}{6.613}{t}
Lat grunnvekta \(\mu_0\) i 3.01 og fordelinga i 6.6 gjelde dei same formene \(D\). Kvart par av agent og omsyn er då eitt trykk, og

\[
\Lambda(x)=\beta\Phi(x)+\lambda_{SG}(x)+c,\quad
\lambda_{SG}(x)=-\log\frac{\mu_{SG}(x)}{\mu_0(x)},
\]

der \(c\) er ein konstant. Landskapet har ein del agentane ber gjennom vurderinga, og ein del grammatikken, stegprosessen og substratet ber utan vurdering (6.003).

Kvar eit trykk høyrer heime, kan prøvast: det som flyttar seg når prioriteringa vert endra, ligg i \(\Phi\); det som står fast, ligg i \(\lambda_{SG}\).
\end{formproposition}

\begin{formproposition}{6.614}{t}
For to former \(x\) og \(y\) i \(D\) er

\[
\log\frac{P(x)}{P(y)}=\log\frac{\mu_{SG}(x)}{\mu_{SG}(y)}-\beta[\Phi(x)-\Phi(y)].
\]

Forholdet avheng ikkje av kva andre former som står i \(D\). Flyttar eit nytt alternativ forholdet mellom to gamle medan grunnvekta og estimata deira står faste, held ikkje modellen med dei oppgjevne ledda.\kjelder{170}
\end{formproposition}

\begin{formproposition}{6.62}{t}
Z \(=\) \(∑_{x∈D}\) \(μ_{\mathrm{SG}}\)(x) exp(\(−\)\(β\) \(Φ\)(x)) normaliserer fordelinga. I det endelege tilfellet med endelege verdiar er 0 \(<\) Z \(<\) \(∞\).

I eit teljeleg rom er Z endeleg når \(μ_{\mathrm{SG}}\) er ei sannsynsfordeling og \(Φ\) er avgrensa nedanfrå; han er positiv når D har positiv grunnvekt.

Ein atomlaus formalgebra (1.611) er ikkje i seg sjølv eit atomlaust sannsynsrom. For eit kontinuerleg operativt rom vert summen erstatta med \(Z=\int_D\exp[-\beta\Phi(x)]\,\mu_{\mathrm{SG}}(dx)\), der grunnmålet må oppgjevast og \(0<Z<\infty\). Når ei ny objektdeling eller akse endrar D, må grunnmålet og koplinga mellom dei to formromma òg oppgjevast.
\end{formproposition}

\begin{formproposition}{6.63}{t}
Grammatikken avgrensar det genererbare, vilkåra avgrensar D og grunnfordelinga, og agentane vektar alternativa gjennom \(Φ\). Haldet vert prøvd gjennom dei faktiske følgjene.

Fastset ledda før nye utfall vert føresagde. Då kan feilen finnast i eit ledd, ikkje løysast ved at alle vert tilpassa utfallet.
\end{formproposition}

\begin{formproposition}{6.631}{t}
Komponentane avgrensar kvarandre. Utan eit steg kan eit føretrekt alternativ verte utilgjengeleg. Utan eit treffande estimat kan eit gjennomført val svikte. Utan kapasitet kan ei rett vurdering verte verknadslaus.

Desse tapa høyrer til ulike ledd og krev ulike endringar.
\end{formproposition}

\begin{formproposition}{6.632}{t}
Fleire agentar kan få plass i same modell, men ein vekta sum er ikkje heile teorien om samspelet. Endrar agentane kvarandre sine mål eller tilgangar undervegs, må desse overgangane inn i forklaringa.
\end{formproposition}

\begin{formproposition}{6.633}{t}
Lat agentane prøve det same forslaget frå ein felles forslagskjerne, med brøken \(r\) frå 6.64. Agent \(i\) godtek uavhengig av dei andre med sannsyn \(\min\{1,\exp(-\beta_i[\hat\Lambda_i(x')-\hat\Lambda_i(x)])\}\), og steget vert gjennomført når alle godtek. Under vilkåra i 6.64 har kjeda då detaljert balanse med

\[
P(x)\propto\mu_{SG}(x)\exp\bigl(-\sum_i\beta_i\hat\Lambda_i(x)\bigr).
\]

Dette er fordelinga i 6.6 med \(\beta=\sum_i\beta_i\) og \(w_i=\beta_i/\beta\). Innverknaden til ein agent er hans del av den samla seleksjonsstyrken.

Avgjer i staden éin tilfeldig vald agent kvart steg åleine, vert den stasjonære fordelinga som regel ei anna. Den vekta summen er difor ein hypotese om korleis agentane avgjer saman, ikkje berre om kor mykje kvar tel.
\end{formproposition}

\begin{formproposition}{6.64}{t}
Lat \(D\) vere endeleg, \(\mu_{SG}>0\) på \(D\), og \(\Phi\) og \(\beta\) vere endelege. Alle ledda står faste under utforskinga. Frå \(x\) føreslår stegprosessen \(x'\) med sannsyn \(q(x,x')\). Forslag utanfor \(D\) vert avviste. For \(q(x,x')>0\) er Metropolis–Hastings-regelen\kjelder{121}

\[
\begin{aligned}
r(x,x')&=\frac{\mu_{SG}(x')q(x',x)}{\mu_{SG}(x)q(x,x')},\\
\alpha(x,x')&=\min\{1,r(x,x')\exp(-\beta[\Phi(x')-\Phi(x)])\}.
\end{aligned}
\]

Brøken er \(r(x,x')\) i den korrigerte politikken (5.7). Når \(q(x',x)=0\), er aksepten null. For \(x'\ne x\) er \(T(x,x')=q(x,x')\alpha(x,x')\); resten av sannsynet vert verande i \(x\). Straumen er

\[
P(x)T(x,x')=\min\{P(x)q(x,x'),P(x')q(x',x)\}.
\]

Han er lik begge vegar, så \(P\) er stasjonær. Dersom dei gjensidig tillatne stega bind heile \(D\) saman og det er positivt sannsyn for å verte ståande i kvar posisjon, konvergerer kjeda til denne eine stasjonære fordelinga. Tida det tek å kome nær henne med oppgjeven presisjon, er blandingstida.\kjelder{110}

Når \(\mu_{SG}(x)q(x,x')=\mu_{SG}(x')q(x',x)\), fell brøken bort og aksepten vert den i 5.7. Ved \(\beta=0\) bevarer den korrigerte prosessen grunnfordelinga.
\end{formproposition}

\begin{formproposition}{6.641}{t}
Når vekter, estimat eller tilgjengelege steg endrar seg før prosessen har blanda seg, er den stasjonære fordelinga ikkje nødvendigvis ei god skildring av utfalla. Det må då forklarast korleis fordelinga utviklar seg i tid.

Eit forslag i ei skisse kan trekkjast tilbake; eit materialsteg kan vere irreversibelt. Når \(q(x',x)=0\), avviser Hastings-regelen steget frå \(x\) til \(x'\). Korreksjonen kan difor ikkje i seg sjølv gje ei likevektsforklaring av slike materielle forløp.

Me må difor skilje den epistemiske søkjekjernen for skisser og simuleringar frå den materielle overgangskjernen for gjennomføringa. Søkjekjernen kan gje fordelinga i 6.6 når vilkåra i 6.64 held og kjeda blandar seg før vilkåra endrar seg. Den materielle kjernen kan vere einsretta og treng ikkje oppfylle detaljert balanse. Ei lokal likevekt i skissa er ikkje ei likevekt i historia.

I eit historisk forløp må tilstanden ta med det minnet dei vidare overgangane krev. Med tidsavhengige kjernar er utviklinga \(p_{t+1}(x')=\sum_x p_t(x)T_t(x,x')\). Fordelinga i 6.6 kan nyttast som eit lokalt valbilete berre der dei oppgjevne vilkåra ber det.
\end{formproposition}

\begin{formproposition}{6.642}{t}
I grensa \(β\) \(→\) \(∞\) vert steg til høgare \(Φ\) avviste. Ein slik lokal prosess kan verte verande i eit lokalt minimum sjølv om den stasjonære grensa i 6.612 legg vekta på dei globale. Ligg alle nabostega minst \(h>0\) høgare, og \(r=1\), er venta opphaldstid i minimumet minst \(e^{\beta h}\) steg.

Lokal rørsle og stasjonær fordeling er difor ulike spørsmål. Historia og dei tilgjengelege stega avgjer om ein dal kan forlatast i tida prosessen har.
\end{formproposition}

\begin{formproposition}{6.65}{t}
Grensetilfella frå 5.72 kan prøvast i fordelingsmodellen:\kjelder{92}

(i) Ved \(β\) \(=\) 0 eller konstant \(Φ\) er P grunnfordelinga avgrensa til D.

(ii) Med éin agent og samla, vekta vurdering \(\hat{Λ}\) som er 0 i K\(+\) og 1 elles, får ei kjend vellukka form \(e^{β}\) gonger så stort sannsyn som ei anna med same grunnvekt. Dette er ei føreseiing om val, ikkje ei utleiing av tanke-kunnskapsrørsla.

(iii) Med fast parameterrom og ei vurdering av alle aksar er modellen ein seleksjonsmodell for optimering. Ein gradientprosess krev fleire vilkår om stega.

(iv) Når \(Φ\) er det faktiske avviket som skal minimerast, samlar den stasjonære fordelinga seg om dei beste tilgjengelege posisjonane når \(β\) \(→\) \(∞\). Det seier ikkje kor raskt dei vert nådde.
\end{formproposition}

\begin{formproposition}{6.66}{t}
Den allmenne modellen bind det genererbare, det vurderte og det gjennomførte til eitt sett med prøver. Han kan forklare meir enn ei rein skildring av formene berre dersom dei faste stega, estimata og vektene gjev betre føreseiingar om nye utfall.

Samanbindinga er eit forskingsprogram (7.17), ikkje ei oppdaga universallov. Ho må tole samanlikning med enklare alternativ.
\end{formproposition}

\begin{formproposition}{6.67}{d}
Den samla effektive prioriteringa langs omsyn k er \(W_{k}\) \(=\) \(∑_{i}\) \(w_{i}\) \(ω_{k}\)(\(A_{i}\)), med \(ω_{k}\) \(=\) 0 utanfor lyskjegla til agenten. \(W_{k}\) måler ei vekt i valprosessen, ikkje kjensle eller gjennomføringskraft. Med omsorg kjem vekta frå andre til (6.7).
\end{formproposition}

\begin{formproposition}{6.671}{t}
Når \(W_{k}\) \(=\) 0, har omsynet inga direkte vekt i \(Φ\). Forma kan likevel verte god eller dårleg langs \(X_{k}\) gjennom \(μ_{\mathrm{SG}}\), V og samanhengen med andre omsyn.

Dersom \(Φ\) berre avheng av koordinatane \(x_{W}\) som estimata med positiv vekt avheng av, får dei andre den vilkårsfordelinga \(μ_{\mathrm{SG}}\) gjev når \(x_{W}\) og D er faste. Dette gjev ikkje uavhengigheit mellom aksane; nettopp vilkårsfordelinga kan bere eit indirekte vern.
\end{formproposition}

\begin{formproposition}{6.6711}{t}
Lat \(T(x)\) vere verdiane av alle estimat med positiv vekt. Då er

\[
P(x\,|\,T)=\mu_{SG}(x\,|\,T).
\]

Innanfor det dei vekta estimata skil, vel grunnfordelinga. Eit mål utan vekt får det ho gjev der.
\end{formproposition}

\begin{formproposition}{6.672}{t}
Med separerbar spenning, flat lokal grunnfordeling og eit estimat som er to gonger deriverbart, med felles krumming ved eit ikkje-degenerert indre minimum, for \(β\) \(>\) 0, \(W_{k}\) \(>\) 0 og \(\hat{Λ}_{k}\)\(″\) \(>\) 0, er

Var(\(x_{k}\)) \(≈\) 1 / (\(β\) \(W_{k}\) \(\hat{Λ}_{k}\)\(″\)).

I fleire dimensjonar, ved eit glatt indre minimum \(x_0\) med positivt definitt Hessian \(H=\nabla^2\Phi(x_0)\) og lokalt flat grunnfordeling, er \(\operatorname{Cov}(x)\approx(\beta H)^{-1}\). Kryssledd i H koplar aksane. Dette er ei lokal Gauss-tilnærming på ein samanhengande skala; ved varierande grunnvekt må krumminga av \(-\log\mu_{\mathrm{SG}}\) òg reknast med.

Spreiinga bestemmer produktet av seleksjonsstyrke, prioritering og krumming. Dei kan skiljast ved å endre eitt ledd medan dei andre står faste. Eit estimat med knekk eller flate i botnen følgjer 6.6722.
\end{formproposition}

\begin{formproposition}{6.6721}{i}
Setehøgda varierer lite, ornamentet mykje. Kropp, produksjon, vane og prioritering kan alle vere årsaka. Spreiinga åleine skil dei ikkje.
\end{formproposition}

\begin{formproposition}{6.6722}{t}
Tilnærminga i 6.672 krev krumming i botnen. Er estimatet avstanden \(|x_k-g_k|\) frå eit målpunkt, har \(\hat\Lambda_k\) ein knekk der. Med same separerbare spenning og flate grunnfordeling, langt frå grensene til \(D\), er

\[
\operatorname{Var}(x_k)=\frac{2\varepsilon_k^2}{(\beta W_k)^2}.
\]

Har målregionen breidd \(L>0\) langs aksen, går \(\operatorname{Var}(x_k)\) mot \(L^2/12\) når \(\beta\to\infty\), ikkje mot null.

Spreiinga følgjer toleransen i begge tilfella. Korleis ho fell med \(\beta W_k\), skil avviksmåla: som \((\beta W_k)^{-1}\) ved kvadratisk mål, som \((\beta W_k)^{-2}\) ved avstand.
\end{formproposition}

\begin{formproposition}{6.673}{t}
For eit endeleg \(D\), skriv \(\beta\Phi(x)=\sum_k\theta_k\hat e_k(x)\), der kvar \(k\) no merkar eit bestemt par av agent og omsyn, \(\hat e_k(x)\) er det toleranseskalerte estimatet \(\hat\Lambda_{ik}(x)\) (6.61), og \(\theta_k=\beta w_i\omega_k(A_i)\) er den samla vekta til dette estimatet. Med omsorg er vekta \(\beta w_i(1-\gamma_i)\omega_k(A_i)\) for eigne omsyn og \(\beta w_Ac_{A\to B,j}\) for biletet av ein annan (5.28). Lat \(D\), \(\mu_{SG}\) og alle \(\hat e_k\) stå faste.

Her er \(\theta_k\) ein fri parameter i den statistiske familien: éin kan endrast medan dei andre står faste. Ei endring av éi normalisert omhugsvekt \(\omega_k(A_i)\) endrar derimot minst éi anna vekt (6.61). Desse to inngrepa må skiljast.

Derivasjon av den normaliserte eksponentialfamilien gjev\kjelder{122}

\[
\frac{\partial P(x)}{\partial\theta_k}=-P(x)\bigl(\hat e_k(x)-\mathbb E[\hat e_k]\bigr).
\]

Difor er \(\partial\mathbb E[\hat e_j]/\partial\theta_k=-\operatorname{Cov}(\hat e_j,\hat e_k)\).

Særleg er \(\partial\mathbb E[\hat e_k]/\partial\theta_k=-\operatorname{Var}(\hat e_k)\leq0\). Meir prioritering kan ikkje auke det estimerte avviket langs same omsyn i middel under desse vilkåra. Forventing og kovarians vert tekne under den same fordelinga \(P\).
\end{formproposition}

\begin{formproposition}{6.6732}{t}
Lat \(D\), \(\mu_{SG}\) og dei toleranseskalerte estimata vere fastsette uavhengig av vala. Då fastset valfordelinga vektene \(\theta\) eintydig dersom og berre dersom kovariansmatrisa \(\operatorname{Cov}(\hat e_j,\hat e_k)\) er positivt definitt. Frå \(\theta\) følgjer \(\beta\), \(w_i\) og \(\omega_k(A_i)\) ved normaliseringa i 6.61, og med omsorg \(\gamma_i\) og omsorgsvektene (5.28).

Utan desse vilkåra fastset vala mindre. Er grunnfordelinga ikkje gjeven, er berre \(\log\mu_{SG}-\beta\Phi\) fastsett. Er toleransen ikkje gjeven, er berre \(\theta_k/\varepsilon_k\) fastsett. Deler to agentar same estimat langs eit omsyn, er berre summen av vektene deira fastsett.
\end{formproposition}

\begin{formproposition}{6.674}{t}
For eit faktisk avvik \(e_j(x)\), målt uavhengig av agenten sitt estimat og halde fast under endringa av prioritering, gjev same utleiing

\[
\frac{\partial\mathbb E[e_j]}{\partial\theta_k}=-\operatorname{Cov}(e_j,\hat e_k).
\]

Når standardavvika er positive, kan den deriverte skrivast \(-\operatorname{corr}(e_j,\hat e_k)\,\sigma(e_j)\sigma(\hat e_k)\), med korrelasjonen mellom faktisk avvik og estimat. Absoluttverdien er høgst \(\sigma(e_j)\sigma(\hat e_k)\) ved Cauchy–Schwarz.

Flyttar ein i staden prioritering frå omsyn m til k innan same normaliserte agentprofil, med \(\beta\), \(w_i\) og alle estimat faste, vert den lokale deriverte \(-\beta w_i\operatorname{Cov}(e_j,\hat\Lambda_{ik}-\hat\Lambda_{im})\). Forteiknet kan då ikkje lesast frå kovariansen med k åleine.

Meir prioritering betrar det faktiske utfallet lokalt når kovariansen er positiv, og gjer det verre når han er negativ. Ved null kovarians er førstederiverten null. For ei endeleg endring må samvariasjonen følgjast langs endringa. Ordnar estimatet formene i same rekkjefølgje som det faktiske avviket, er kovariansen ikkje-negativ under kvar fordeling på \(D\), og betringa held for kvar endeleg auke av ei fri vekt.

Er følgja uviss også for ei gjeven form, gjeld identiteten med middelverdien av \(e_j\) for kvar form, så lenge denne ikkje avheng av prioriteringa.

Omhug og kunnskap er difor ulike verkemåtar også i likninga: prioriteringa endrar styrken, estimatet avgjer kva retning styrken tener. Endrar prioriteringa òg modellen, følgjene eller dei tilgjengelege alternativa, gjeld ikkje denne partielle deriverte åleine.
\end{formproposition}

\begin{formproposition}{6.675}{t}
Auka prioritering av eitt omsyn kan verne eller skade eit anna gjennom samvariasjonen i 6.674. Det gjeld også omsyn utan eiga vekt.

Eit indirekte vern er difor eit føreseieleg samband mellom omsyn, ikkje eit unntak frå at omhug verkar gjennom bestemte steg.
\end{formproposition}

\begin{formproposition}{6.68}{t}
Forma vert vald gjennom vurderingar og gjennomført i vilkår vurderingane ikkje rår over. Dei to kan samsvare eller skilje lag. Skilnaden er noko læra skal føreseie og prøve.
\end{formproposition}

\begin{formproposition}{6.7}{t}
Med omsorg får måla til andre vekt i forma. Målvekta, den samla vekta på målet til \(B\) langs \(X_j\), er

\[
W_{B,j}=w_B(1-\gamma_B)\omega_j(B)+\sum_{A\ne B}w_Ac_{A\to B,j}.
\]

Når ingen lyskjegle er tom, er \(\sum_{B,j}W_{B,j}=1\). Omsorg flyttar vekt. Ho lagar ikkje vekt.

Den som lever med forma utan innverknad på henne, har \(w_B=0\). Målet hans har då vekt berre gjennom omsorga til andre, og berre slik dei ser det.
\end{formproposition}

\begin{formproposition}{6.701}{t}
Lat \(x\) og \(x'\) vere like i alle andre vekta estimat. Då er

\[
\begin{aligned}
\log\frac{P(x)}{P(x')}={}&\log\frac{\mu_{SG}(x)}{\mu_{SG}(x')}\\
&-\theta\bigl[\hat\Lambda_{A|B,j}(x)-\hat\Lambda_{A|B,j}(x')\bigr],
\end{aligned}
\]

med \(\theta=\beta w_Ac_{A\to B,j}\). Omsorg kan difor målast i val mellom former som berre skil seg i biletet av avviket hos den andre.

Ved bruk (5.282) fell skilnaden bort der avviket hos den andre ikkje lenger tener agenten sine eigne mål.
\end{formproposition}

\begin{formproposition}{6.71}{t}
Omsorga verkar gjennom biletet. Lat \(\theta=\beta w_Ac_{A\to B,j}\) vere den frie vekta til biletet. Under vilkåra i 6.674 er

\[
\frac{\partial\mathbb E[e_{B,j}]}{\partial\theta}=-\operatorname{Cov}\bigl(e_{B,j},\hat\Lambda_{A|B,j}\bigr).
\]

Omsorg minskar det faktiske avviket hos den andre når biletet følgjer avviket, og aukar det når biletet er vendt feil. Styrken i omsorga avgjer kor mykje. Biletet avgjer kva veg.

Flyttar \(A\) vekt frå eigne omsyn til omsorga, er den deriverte \(-\beta w_A\operatorname{Cov}(e_{B,j},\hat\Lambda_{A|B,j}-\Phi^0_A)\), der \(\Phi^0_A=\sum_k\omega_k(A)\hat\Lambda_{A,k}\) er vurderinga av eigne omsyn. Det \(A\) gjev opp, kan sjølv ha verna \(B\) (6.675).
\end{formproposition}

\begin{formproposition}{6.711}{t}
Lat \(\hat e\) vere biletet og \(e\) det faktiske avviket hos den andre, begge i toleransen hans, og lat berre vekta \(\theta\) til biletet endre seg. Svarkurva

\[
h(u)=\mathbb E[e\,|\,\hat e=u]
\]

er det faktiske avviket i middel mellom formene der biletet har verdien \(u\). Innanfor kvar verdi av biletet står fordelinga fast (6.6711). Svarkurva avheng difor ikkje av \(\theta\), og

\[
\frac{\partial\mathbb E[e]}{\partial\theta}=-\operatorname{Cov}\bigl(h(\hat e),\hat e\bigr).
\]

Fell svarkurva ingen stad, kan meir omsorg ikkje auke det faktiske avviket hos den andre i middel, same kor stor endringa er. Stig ho ingen stad, kan meir omsorg ikkje minske det. Er ho flat, endrar omsorga ingenting for han.

Omsorga treng ikkje måle avviket til den andre rett. Ho må ordne alternativa rett.
\end{formproposition}

\begin{formproposition}{6.712}{t}
Ordninga kan ikkje sikrast innanfrå. Målregionen er den andre sin, biletet er agenten sitt (5.28), og agenten har ikkje noko ukoda møte å samanlikne biletet med (5.0251). Berre eit svar frå den som har avviket, eller ei uavhengig prøve av følgjene (6.674), kan vise svarkurva.

I rekninga svarar det å lukke sløyfa (7.25) til å setje \(h(\hat e)\) i staden for \(\hat e\). Då er \(h(u)=u\), den deriverte er \(-\operatorname{Var}(\hat e)\), og \(\mathbb E[e]=\mathbb E[\hat e]\) for kvar \(\theta\). Éi retting held så lenge berre omsorgsvekta endrar seg. Endrar andre vekter, estimat eller steg seg, må svaret hentast på nytt. Svaret kan òg endre målet (7.243).
\end{formproposition}

\begin{formproposition}{6.7121}{i}
Designaren ventar minst smerte på det mjuke setet (7.12). Etter ei heil vakt er det der folk har mest vondt: svarkurva fell. Før prøvesitjinga gjer meir omsorg skaden større. Set ho svaret i staden for biletet, gjer den same omsorga han mindre.
\end{formproposition}

\begin{formproposition}{6.713}{t}
Lat \(Y(x)\) vere svaret den andre gjev på forma \(x\), og lat biletet byggje på svaret, \(\hat e=f(Y)\). Då er

\[
\begin{aligned}
|\operatorname{Cov}(e,\hat e)|&\leq\eta_Y\,\sigma(e)\,\sigma(\hat e),\\
\eta_Y^2&=\frac{\operatorname{Var}(\mathbb E[e\,|\,Y])}{\operatorname{Var}(e)}.
\end{aligned}
\]

Biletet kan ikkje følgje den andre finare enn svaret skil. Eit finare svar forklarer aldri mindre. Legg svaret noko til det modellen til designaren alt forklarer, kan det ikkje erstattast av modellen (7.201).

For den som ikkje kan svare sjølv, må ein annan kanal bere svaret. Utan kanal er \(\eta_Y=0\), og retninga på omsorga er uprøvd.
\end{formproposition}

\begin{formproposition}{6.72}{t}
Lat \(D\) vere endeleg. Når \(\beta\to\infty\), samlar massen seg på minimum av \(\Phi\) (6.612). Har alle estimata og bileta positiv vekt, er kvart slikt minimum Pareto-optimalt for dei: inga form i \(D\) er minst like god langs alle og betre langs eitt.\kjelder{115}

Har eit mål null vekt, kan eit minimum vere dominert, men berre langs mål utan vekt. Tapet der er ingen annan sin gevinst.

Ei lita, positiv omsorgsvekt vel då, mellom formene som alt var best for dei andre, dei som er best for den som mangla vekt. Ho endrar ingenting anna. Skade som ingen vinn på, står berre att der målet er utan vekt.

Optimaliteten gjeld estimata. For dei faktiske avvika gjeld ho berre så langt estimata ordnar formene rett.
\end{formproposition}

\begin{formproposition}{6.73}{t}
Omsorga held fram gjennom den ho gjeld. Tek \(A\) inn heile vurderinga til \(B\), med bileta og vektene til \(B\), også det \(B\) sjølv har omsorg for, er

\[
\Phi_A=(1-\gamma_A)\Phi^0_A+\sum_{B\ne A}c_{AB}\Phi_B,
\]

med \(\gamma_A=\sum_Bc_{AB}\leq1\). Dette er grensetilfellet der \(A\) kjenner vurderinga til \(B\) utan feil; med eit bilete av henne legg feila seg saman langs vegen.

Følg omsorga frå \(A\): vidare til \(B\) med sannsyn \(c_{AB}\), og bli ståande hos \(B\) med sannsyn \(1-\gamma_B\). Når kvar agent anten har eigne omsyn, \(\gamma<1\), eller ein omsorgsveg til ein som har det, endar kvar slik vandring, og likningane har éi løysing:\kjelder{110,171}

\[
\Phi_A=\sum_B\chi_{AB}\Phi^0_B.
\]

\(\chi_{AB}\) er sannsynet for at vandringa frå \(A\) endar hos \(B\). I \(\Phi=\sum_Aw_A\Phi_A\) (6.6) tel \(B\) då med innverknaden \(\tilde w_B=\sum_Aw_A\chi_{AB}\).
\end{formproposition}

\begin{formproposition}{6.731}{t}
Ein lukka krins der ingen har eigne omsyn, fastset inga vurdering: kvar funksjon som er felles for heile krinsen, oppfyller likningane i 6.73.

To som berre vil det den andre vil, vil ingenting. Omsorga må ende i eit omsyn som er nokon sitt eige.
\end{formproposition}

\begin{formproposition}{6.74}{t}
Lat \(Q\) vere ei samling agentar med positiv innverknad, og lat \(\beta>0\). Valfordelinga er uendra under kvar omfordeling av innverknaden innanfor \(Q\) dersom og berre dersom vurderingane \(\Phi_A\), \(A\in Q\), er like på \(D\) opp til konstantar.

Då ber vurderinga ikkje merke om kven avviket høyrer til (6.121), og samlinga vel som éin agent. Det spelar inga rolle kven av dei som avgjer.

Dette gjev ei prøve: endrar forma seg med kven i samlinga som rår, er spenninga ikkje delt langs dei omsyna der endringa viser seg. Lik vurdering kan òg kome av like mål utan kopling. Prøva kan difor avkrefte delt spenning, ikkje stadfeste henne.
\end{formproposition}

\begin{formproposition}{6.741}{t}
I ein samanhengande krins, lat eigenandelane \(1-\gamma_A\) gå mot null i fast innbyrdes forhold \(s_A\). Då går \(\chi_{AB}\) for kvar \(A\) mot

\[
\frac{\pi_Bs_B}{\sum_{B'}\pi_{B'}s_{B'}},
\]

der \(\pi\) er den stasjonære fordelinga til omsorgsvandringa. Vurderingane fell saman, og innverknaden innanfor krinsen sluttar å telje (6.74).

Delt spenning (6.121) er grensa for gjensidig omsorg. Grensa vert ikkje nådd: i sjølve grensa fastset krinsen inga vurdering (6.731).
\end{formproposition}

\begin{formproposition}{6.742}{t}
Lat omsorga vere sterk innanfor grupper og svak mellom dei, med små eigenandelar. Kvar gruppe \(Q\) vurderer då tilnærma som éin agent (6.741), og gruppene har omsorg for kvarandre med vektene

\[
\bar c_{QQ'}=\frac{u_{QQ'}}{\bar s_Q+\sum_{Q''}u_{QQ''}}.
\]

Her er \(\bar s_Q=\sum_{A\in Q}\pi_A(1-\gamma_A)\) eigenandelen og \(u_{QQ'}=\sum_{A\in Q}\pi_A\sum_{B\in Q'}c_{AB}\) omsorga ut, begge vekta med den stasjonære fordelinga i gruppa. Likningane i 6.73 gjeld då tilnærma på neste skala.\kjelder{119,172}

Omsorga har same form på kvar skala. På neste skala tel ikkje kor sterk ho er innanfor gruppa, berre forholdet mellom det gruppa gjev ut og det ho held for seg sjølv.

Dette gjeld vurderinga. Ein akse ingen del har i lyskjegla si (6.123), kjem ikkje av omsorg mellom delane åleine.
\end{formproposition}

\begin{formproposition}{6.75}{t}
Lat \(Q\) vere ei samling som er éin agent (6.123) med ei hinne om alle (6.211). Når ein del svarar på avviket hos ein annan del langs eit omsyn heilskapen er avhengig av, og svaret minskar avviket, er det regulering innanfor hinna som motverkar ei forstyrring som elles ville løyse henne opp (5.021).

Under dette vilkåret er omsorg mellom delane sjølvhald i heilskapen (5.115): eitt svar, skildra på to skalaer.

Vilkåret er avhengigheita. Omsorg for ein som heilskapen ikkje er avhengig av, er framleis omsorg. Ho er ikkje sjølvhald i denne heilskapen.\kjelder{22,23,87}
\end{formproposition}

\begin{formproposition}{6.751}{t}
Det omvende gjeld ikkje. Ein heilskap kan halde seg utan omsorg mellom delane: éin del kan styre vilkåra dei andre svarar på (6.111), bruke avvika deira for eigne mål (6.021), eller føre gevinsten ein del tek på det felles, tilbake til han som avvik (6.1213).

At ein heilskap held seg, viser difor ikkje at delane har omsorg for kvarandre. Det må prøvast hos delane, ved arbeidet avviket hos den andre set i gang (5.281).
\end{formproposition}

\begin{formproposition}{6.752}{o}
I froskeembryo får celler som eit kreftgen driv mot svulstliknande vekst, eit avvikande membranpotensial før veksten er synleg. Vert potensialet tvinga tilbake, vert det færre slike danningar, endå genet framleis er uttrykt.\kjelder{174}

I mange svulstar er dei direkte kanalane mellom cellene svekte. Sambandet er ikkje eintydig: det skiftar med svulsttype og stadium, og i eitt forsøk i det same dyret gav broten kopling færre svulstar og sterkare kopling fleire.\kjelder{175,176}

Det som er vist, er at tilstanden i vevet er med og avgjer om cella held seg til forma. Det er ikkje vist at svulstar allment kjem av tapt kopling.
\end{formproposition}

\begin{formproposition}{6.753}{a}
Ei lesing er at svulstcella har fått eit mindre målrom. Ho regulerer framleis kompetent, men avvika i vevet set ikkje lenger arbeid i gang i henne: lyskjegla har krympa til cella sjølv (5.26).\kjelder{22,85,173}

Lesinga er ein hypotese i biologien. Ho skil seg frå tapt evne ved det ho føreseier: cella held sitt eige like godt eller betre (6.754), og vevet kan ta opp att forma når koplinga vert retta, utan at cella vert bytt ut. Der svulsten held fram med retta kopling, held ikkje lesinga for det tilfellet.
\end{formproposition}

\begin{formproposition}{6.754}{t}
Skriv vurderinga til \(A\) som \(\Phi_A=(1-\gamma_A)\Phi^0_A+\gamma_A\Psi_A\), der \(\Phi^0_A\) gjeld eigne omsyn og \(\Psi_A\) er bileta av avvika hos dei andre, vekta innbyrdes som før. Lat berre omsorgsdelen \(\gamma_A\) endre seg. Då er (6.674)

\[
\frac{\partial}{\partial\gamma_A}\mathbb E\bigl[\Psi_A-\Phi^0_A\bigr]=-\beta w_A\operatorname{Var}\bigl(\Psi_A-\Phi^0_A\bigr)\leq0.
\]

Når omsorgsdelen fell, veks avstanden mellom avviket agenten ser hos dei andre, og avviket langs eigne omsyn. Der dei to ikkje samvarierer positivt, fell det eigne avviket medan avviket hos dei andre stig.

Dette er signaturen på eit krympa målrom: agenten held sitt eige like godt eller betre. Tapt evne gjev ein annan signatur: då sviktar òg det agenten sjølv prioriterer.
\end{formproposition}

\begin{formproposition}{6.755}{a}
Den same krympinga skal kunne finnast i verkstader, føretak og fag. Vert kanalen som bar avvika mellom delane, broten medan delane elles står uendra, skal svikten auke i det dei heldt saman: i overgangane mellom delane og langs omsyn ingen del har som sitt. Det kvar del held åleine, skal stå like godt eller betre (6.754).

Fastset før brotet kva omsyn som er eigne og kva som er felles. Mål sviktrisikoen (7.152) langs begge, mot ei samling utan brot.

Hypotesen fell dersom svikten aukar like mykje langs dei eigne omsyna; då er det evna som er tapt. Ho fell dersom det felles haldet står uendra. Og ho fell dersom haldet ikkje kjem tilbake når kanalen vert sett inn att med dei same delane.

Hypotesen gjeld kanalen som ber avvika, ikkje mengda samband. Sterkare kopling er ikkje alltid betre (6.122), og ei kopling kan òg bere eit hald som skadar dei utanfor (7.263).
\end{formproposition}

\begin{formproposition}{6.76}{t}
Vekta når lenger enn omsorga til agenten: gjennom dei han har omsorg for, tel òg dei som dei har omsorg for. Målet til \(B\) har vekt i forma viss og berre viss han sjølv har innverknad, eller det går ein veg av omsorg til han frå nokon som har (6.7, 6.73).
\end{formproposition}

\begin{formproposition}{7}{t}
Forma viser kor langt omhugen rakk.
\end{formproposition}

\begin{formproposition}{7.001}{t}
Omhugen forma viser, er den som vart verksam: der vekt i vurderinga (6.61), eit estimat som ordna formene rett (6.674), og kapasitet (5.21) møttest i eit hald (5.231). Det gjeld omhug for eigne mål og omsorg for måla til andre (5.28).
\end{formproposition}

\begin{formproposition}{7.01}{t}
Forma verkar med heile ordninga ho har, også der agentane berre held nokre relasjonar i modellen. Følgjene stoggar ikkje ved grensa for omsyna deira.
\end{formproposition}

\begin{formproposition}{7.011}{d}
Verknadskjegla til ei form er alt valet av henne gjer skilnad for: kvar agent, kvart omsyn og kvart tidsrom der avviket hadde vore eit anna under minst eitt av alternativa som stod opne (6.6).

Kjegla er gjeven ved følgjene, ikkje ved kva nokon veit om dei (1.2453). Det kjende utsnittet veks med prøvene (1.52).

Ein skilnad er vesentleg når han er større enn ein del av toleransen til den som ber han, fastsett før prøva (5.11, 7.151).
\end{formproposition}

\begin{formproposition}{7.012}{t}
Lyskjegla er det agenten har som mål (5.22). Verknadskjegla er det forma verkar på. Ingen av dei set grensa for den andre (5.222, 6.027, 7.01).

Det som ligg i verknadskjegla og utanfor lyskjegla til alle med innverknad, har inga vekt i valet (6.7). Det får vekt på to måtar: ved at den som ber følgja, får innverknad, eller ved at nokon med innverknad har omsorg for han.
\end{formproposition}

\begin{formproposition}{7.013}{d}
Det som ligg i verknadskjegla utan å vere halde av nokon (5.231), heller ikkje av den som ber følgja, er det uhaldne (5.2321). Det uverna er den delen av det uhaldne som heller ikkje får vern gjennom eit anna omsyn (6.675) eller gjennom grunnfordelinga (6.671).

Der avgjer substratet og variasjonen i stega utfallet (7.02), og eit omsyn med vekt kan skade det. Vernet gjennom eit anna omsyn er eit samband, ikkje eit hald: det kan snu når vilkåra endrar seg.
\end{formproposition}

\begin{formproposition}{7.02}{t}
Langs eit omsyn kan utfallet vere halde gjennom eige arbeid, verna gjennom eit anna omsyn, gjeve av substratet eller overlate til variasjonen i stega. Desse vegane kan ende i same form.

Forma viser omhugen først når ho vert lesen saman med forløpet: det som skjer når vilkår eller prioriteringar vert endra.
\end{formproposition}

\begin{formproposition}{7.021}{t}
Omhug utan kapasitet kan setje sterke spor utan å lukkast. Kapasitet utan omhug kan verte liggjande ubrukt. Forma viser kva som vart realisert av arbeidet; for å skilje det meinte frå det moglege og det haldne trengst ulike prøver.
\end{formproposition}

\begin{formproposition}{7.03}{t}
Der eit omsyn førte til arbeid som heldt forma under forstyrring, er haldet eit spor etter omhugen: ein kausal samanheng mellom registrert avvik, svar og utfall.
\end{formproposition}

\begin{formproposition}{7.031}{t}
Materie er oppleving sett utanfrå. Det levande er eit sjølv sett utanfrå. Langs dei omsyna der målretta verksemd bar forma, er forma sporet av omhugen, sett utanfrå.

Forma viser omhugen som vart verksam, ikkje heile innsida til den som forma.
\end{formproposition}

\begin{formproposition}{7.04}{t}
Omhugen ein gjenstand viser, er lånt: arbeidet til dei som forma og brukte han, festa i materialet. Berre der gjenstanden sjølv har reguleringa i 5.11, har han eigen omhug. Målregionen kan også då vere sett av andre (5.116).
\end{formproposition}

\begin{formproposition}{7.041}{t}
Gjenstanden er eit arkiv i den grad tidlegare arbeid held fram som verknad i han. Arkivet kan bevare eit omsyn ingen no arbeider med, og miste eit omsyn nokon framleis prøver å halde.
\end{formproposition}

\begin{formproposition}{7.1}{a}
Den realiserte forma ber ein sviktsignatur der utelatingar i arbeidet fører til systematiske feil under dei trykka som vart utelatne.

Avgrensa hald gjev bestemte sviktmønster. Prøva krev at grensa for haldet vert fastsett før dei nye sviktane vert lesne.

Blindskap, likegyldigheit og avmakt kan gje liknande utfall. Det same kan konflikt mellom omsyn som kvar for seg kan haldast (5.231), og eit estimat som peikar feil (6.674). For å skilje dei må kunnskap, prioritering, kapasitet og estimat undersøkjast kvar for seg.
\end{formproposition}

\begin{formproposition}{7.11}{t}
Ei lesing av signaturen bind sporet til steg, vilkår og prioriteringar. Ho vert prøvd ved at ei endring av eitt ledd endrar nye utfall slik forklaringa sa.
\end{formproposition}

\begin{formproposition}{7.12}{a}
Når alternativ, vilkår og handlag står faste, skal endra prioritering endre nye val slik 6.673 føreseier. Samsvaret mellom estimat og faktisk avvik avgjer kva det gjer med forma i bruk (6.674).

Lat ein designar velje mellom to sete. Ho ventar at den som skal sitje ei heil vakt, får minst vondt på det mjuke. Jo meir smerta tel i valet, jo oftare vil ho velje det mjuke setet. Viser prøva at det er der folk får mest vondt, får den same feilvurderinga større følgjer jo sterkare omsorga hennar er.
\end{formproposition}

\begin{formproposition}{7.121}{a}
Mål biletet til designaren og det faktiske avviket hos den andre på dei same formene, før valet. Svarkurva (6.711) føreseier då forteiknet: aukar omsorgsvekta medan alt anna står fast, skal det faktiske avviket falle i middel når kurva stig, og stige når ho fell.

Skjer det motsette, held ikkje valmodellen for denne prøva.
\end{formproposition}

\begin{formproposition}{7.13}{t}
Arkivet er eit grunnlag for gransking, ikkje ein ferdig dom. Ei kausal lesing må skilje mellom fleire forløp som kunne ha gjeve same spor.\kjelder{72}
\end{formproposition}

\begin{formproposition}{7.14}{d}
Sviktlatensen langs \(X_k\) er \(\rho_k(A)=\Delta_k/\tau^{\mathrm{pred}}_k(A)\): tida \(\Delta_k\) før konsekvensen viser seg, delt på prediksjonshorisonten (5.114), når denne er positiv. Begge vert rekna frå same starttid.

Tre horisontar må skiljast: \(\tau^{\mathrm{pred}}_k\) for kva agenten kan føreseie med oppgjeven presisjon, \(\tau^\Gamma_k\) for kor langt målet rekk fram i tid (5.22), og \(\tau_k\) for kva ein tilgjengeleg politikk kan halde under prøvevilkåra (5.114). Sviktlatensen måler berre forholdet til den første.
\end{formproposition}

\begin{formproposition}{7.141}{t}
Ved \(ρ\) \(≤\) 1 ligg konsekvensen innanfor den prøvde prediksjonshorisonten; ved \(ρ\) \(>\) 1 utanfor. Ved null horisont og positiv \(Δ\) set eg \(ρ\) \(=\) \(∞\); ved \(Δ\) \(=\) 0 er \(ρ\) \(=\) 0.

Eit utfall utanfor horisonten kan vere kjent som risiko. Nye prøver eller ei vidare representering kan flytte grensa.
\end{formproposition}

\begin{formproposition}{7.142}{i}
Ei prøvesitjing som berre ber føreseiingar éin time fram, møter smerte etter åtte timar med \(ρ\) \(=\) 8. Ei lengre prøve kan flytte horisonten.

Bruker nokon ein kopp i tjue minutt, er dette brukstida, ikkje nødvendigvis prediksjonshorisonten. Den som kjenner materialet, kan føre følgjer hundreår fram inn i valet no.
\end{formproposition}

\begin{formproposition}{7.143}{t}
Ved \(\rho_k\leq1\) ligg følgja innanfor prediksjonshorisonten og kan føreseiast med oppgjeven presisjon. Er \(\Delta_k>\tau^\Gamma_k(A)\), fell følgja utanfor målhorisonten i lyskjegla. Med eit tilgjengeleg vern er dette likegyldigheit i tid; utan både mål og kapasitet er det blindskap i tydinga i 5.232. Er \(\Delta_k\leq\tau^\Gamma_k(A)\), men \(\Delta_k>\tau_k(A)\), gjev det avmakt i tid (5.2322).

Er \(\Delta_k\leq\min(\tau^\Gamma_k(A),\tau_k(A))\), har agenten både mål og eit vern tilgjengeleg langs omsynet under prøvevilkåra. Kva som vert halde, kjem då an på politikken han tek i bruk; vern som finst kvar for seg, treng ikkje kunne nyttast samstundes (5.231).

Ved \(\rho_k>1\) ligg følgja utanfor den prøvde prediksjonshorisonten. Lokal tilbakekopling eller eit vern sett i verk no kan likevel halde omsynet når følgja kjem. Blindskap kan finnast på begge sider av prediksjonsgrensa; han gjeld kva agenten bryr seg om og kan halde (5.232).
\end{formproposition}

\begin{formproposition}{7.144}{t}
Sviktlatensen er relativ til den som undersøkjer. Det som ligg utanfor prediksjonshorisonten til den som formgav, kan liggje innanfor ein annan sin.

Den andre kan òg ha eit anna mål og ein annan toleranse. Kva som tel som svikt, må difor følgje med når prøver vert samanlikna.
\end{formproposition}

\begin{formproposition}{7.15}{t}
Ei forklaring av sviktsignaturen må skilje omsyn, gjennomføring og følgjer. Omhug, lyskjegle og kapasitet gjer skilja til spørsmål som kan prøvast.\kjelder{21,70}
\end{formproposition}

\begin{formproposition}{7.151}{d}
Ei førehandsprøve av læra byrjar før dei nye utfalla er kjende.

Fastset formrom, aksar, tillatne steg, grunnfordeling, mål, toleransar og forstyrringar. Registrer modellen sine føreseiingar om kva som vil skje, og agenten sine eigne estimat før valet. Prøv prioritering og kapasitet kvar for seg.

Hald deretter nye tilfelle utanfor tilpassinga. Samanlikn dei føresagde vala og følgjene med det som faktisk skjer, og med ein enklare grunnmodell under same vilkår. Ei rett føresagd handling med feil følgje er ikkje same feilen som ei rett vurdering agenten ikkje kunne gjennomføre.
\end{formproposition}

\begin{formproposition}{7.152}{d}
Lat \(e_{k}\)(x) vere faktisk avvik for forma \(x\) i ei uavhengig konsekvensprøve (5.11), og \(ε_{k}\) den fastsette toleransen. Den føresagde sviktrisikoen er\kjelder{120}

\(R_{k}\) \(=\) \(∑_{x∈D}\) P(x) \(𝟙\)[\(e_{k}\)(x) \(>\) \(ε_{k}\)].

Valmodellen gjev P; konsekvensprøva gjev \(e_{k}\). Begge trengst for å føreseie skade. Er følgja uviss også når forma er gjeven, står sannsynet for svikt i staden for indikatoren.
\end{formproposition}

\begin{formproposition}{7.1521}{t}
Ved faste vilkår som i 6.674, med fast toleranse, er

\[
\frac{\partial R_j}{\partial\theta_k}=-\operatorname{Cov}\bigl(\mathbf 1[e_j>\varepsilon_j],\hat e_k\bigr).
\]

Her merkar \(j\) følgjeaksen og \(k\) eit par av agent og omsyn frå 6.673. Når \(\hat e_k\) ikkje minkar der \(e_j\) veks, er kovariansen ikkje-negativ: sviktindikatoren minkar ikkje når avviket veks. Estimatet treng då berre ordne formene rett, ikkje treffe storleiken. Meir fri prioritering av eit slikt estimat kan ikkje auke sviktrisikoen i modellen, heller ikkje ved ei endeleg endring. Eit feilvendt estimat kan snu retninga.
\end{formproposition}

\begin{formproposition}{7.153}{t}
Når dei faste vilkåra i 6.673 er oppfylte, skal auka fri prioritering av eit estimert avvik ikkje auke det estimerte avviket i middel. Skjer det likevel, held ikkje valmodellen for denne prøva.

Prøva gjeld meir enn forteiknet. Kovariansen målt før inngrepet skal føreseie kor mykje middelverdien flyttar seg ved ei lita endring (6.674). Agentar med ulike vekter må prøvast kvar for seg: i ei blanding rommar den samla kovariansen òg skilnaden mellom agentane.

Vert valet føresagt rett, men skaden feil, må konsekvensmodellen eller dei endra vilkåra granskast. Fell ein oppgjeven kapasitet under avtalte forstyrringar, held ikkje kapasitetsmålet. Feilar modellen systematisk mot ein enklare, førehandsvald grunnmodell på nye tilfelle, har den ekstra teorien enno ikkje vist forklaringsarbeidet sitt.

Ei lære er ikkje prøvbar dersom kvart motdøme kan gjerast til ei stadfesting i ettertid.
\end{formproposition}

\begin{formproposition}{7.16}{t}
Ei lære som gjer eitt ufullstendig omsyn til einaste mål, gjev ei tydeleg prøve: fastset utelatinga, føresei følgjene hennar, og prøv dei. Slik vert uskuldig forenkling skild frå forenkling som systematisk skadar det ho ikkje held.
\end{formproposition}

\begin{formproposition}{7.17}{d}
Læra er eit forskingsprogram (6.66): ein hard kjerne som vert halden fast, hjelpehypotesar som vert prøvde og kan skiftast ut, føreseiingar som er fastsette før utfalla, og ein regel for kva som tel som framgang.\kjelder{165}
\end{formproposition}

\begin{formproposition}{7.171}{a}
Den harde kjernen er arkitekturen i 5.724 og påstandane som ber han.

Å forklare ei form er å forklare utvalet mellom alternativ (1.13). Generering, vekting og gjennomføring er ulike ledd, og saman strekkjer dei til (5.7231). Prioritering, estimat og kapasitet er ulike årsaker: kvar kan endrast medan dei andre står faste, og kvar sviktar på sin måte (5.201, 5.232). Omhug for eigne mål og omsorg for måla til andre er ulike vekter (5.28). Val og følgje krev kvar si prøve (6.611). Alt dette gjeld same kva agenten er laga av, og på kvar skala (5.111, 6.11).

Kjernen er sjølv ein hypotese. Han fell etter regelen i 7.173, ikkje ved eitt motdøme. Postulata om oppleving (1.23, 1.24) og forpliktinga i 7.26 vert ikkje prøvde av programmet, og inga føreseiing i 7.174 kviler på dei (1.2342).
\end{formproposition}

\begin{formproposition}{7.172}{d}
Hjelpehypotesane gjev den harde kjernen eit innhald som kan prøvast: valet av formrom og aksar (1.321), additiviteten i landskapet (3.01), kjernelikninga for stega (5.7), grunnfordelinga og den eksponentielle valmodellen (6.6), og kvar måte å måle omhug, omsorg, kapasitet og avvik på. Følgjene i 6.673–6.76 står og fell med valmodellen.

Kvar hjelpehypotese kan fellast utan at kjernen fell. Ingen kan skiftast ut utan at skiftet vert ført opp (7.173).
\end{formproposition}

\begin{formproposition}{7.173}{d}
Eit skifte i hjelpehypotesane er framgang når det føreseier eit nytt utfall som så vert funne (7.151), og læra framleis føreseier betre enn den enklare grunnmodellen. Eit skifte som berre bergar den harde kjernen frå eit motdøme, er eit tap og vert ført opp som tap (7.153).

Kjernen fell når skifta berre bergar. Kor mange tap på rad han toler, vert fastsett før rekkja byrjar, ikkje etter.
\end{formproposition}

\begin{formproposition}{7.174}{a}
Programmet står på desse føreseiingane. Kvar vert fastsett før utfallet (7.151) og fell når det motsette vert funne.\kjelder{166}

(i) Forholdet mellom sannsyna for to former endrar seg ikkje når andre alternativ kjem til eller fell bort (6.614).

(ii) Når éi fri vekt aukar litt, flyttar kvart estimert avvik seg i middel så mykje som kovariansen målt før endringa seier, agent for agent (6.673, 7.153).

(iii) Forteiknet på kovariansen mellom estimat og faktisk avvik, målt før valet, føreseier om meir prioritering minskar eller aukar den faktiske skaden (6.674, 7.12).

(iv) Spreiinga langs ein akse fell med prioriteringa etter ein eksponent som følgjer av avviksmålet og er oppgjeven før prøva (6.672, 6.6722).

(v) Sannsynet for å godta eit steg oppover fell eksponentielt med den estimerte skilnaden og avheng ikkje av vegen fram til steget (5.7).

(vi) Svikt kjem oftare langs aksar som på førehand er fastsette som uhaldne og uverna, enn langs haldne, og så ofte som \(R_k\) seier (7.1, 7.152).

(vii) Blindskap, likegyldigheit og avmakt vert retta av ulike inngrep: meir kapasitet rettar ikkje likegyldigheit, og meir prioritering rettar ikkje avmakt (5.232, 5.271).

(viii) Svarkurva flyttar seg ikkje når omsorgsvekta vert endra. Stig ho, minskar meir omsorg det faktiske avviket hos den andre (6.711, 7.121).

(ix) Der spenninga er delt, endrar ikkje forma seg med kven som avgjer. Der koplinga brest, løyser delane framleis sine eigne oppgåver medan det felles haldet fell (6.74, 6.755).

(x) Når steg, vurdering og hald er fastsette kvar for seg, gjev ikkje eit fjerde ledd betre føreseiingar om nye tilfelle (5.7231).

(xi) Stilkategoriar som er fastsette utan å sjå på formmåla, føreseier nye former betre enn materialet, med like mange parametrar (2.202).

(xii) Tradisjonar utan overføring, men under same vilkår, ligg nærare kvarandre enn tilfeldig vandring frå same start gjev (3.41).

Ei føreseiing kan strykast frå lista berre ved å førast opp som tap (7.173).
\end{formproposition}

\begin{formproposition}{7.2}{t}
Designhandlinga endrar vilkåra ein annan agent navigerer under. Ho kan endre kva han kan gjere, kva han må tole, og kva han lærer å vente.

Ei vitskapleg undersøking av handlinga må følgje denne overføringa. Det ferdige produktet er ikkje slutten på prosessen som skal forklarast.
\end{formproposition}

\begin{formproposition}{7.201}{t}
Ei endring av vilkåra til eit sjølv er innanfrå ei endring av det det møter og kjenner. Svaret frå den som ber følgjene, høyrer til kunnskapen om forma. Der det legg noko til modellen designaren har av svaret, kan det ikkje erstattast av modellen (6.713).
\end{formproposition}

\begin{formproposition}{7.21}{t}
Eit avvik som ingen tok inn i arbeidet, kan verte liggjande hos den som møter forma. Også ein svikt som vart føresagd og forsøkt hindra, kan verte lagd over på andre.

Fordelinga av innsikt, avgjerdsmakt og kostnad må difor undersøkjast kvar for seg.
\end{formproposition}

\begin{formproposition}{7.22}{t}
Der eit sviktmønster rammar eit menneske, kan det som utanfrå er ein systematisk feil, innanfrå vere smerte, tap eller avmakt. Ei forklaring som finn mønsteret, har ikkje dermed gjort greie for kva det er å bere det.
\end{formproposition}

\begin{formproposition}{7.23}{d}
Det predikerte sviktrommet er omsyna der \(R_{k}\) ligg over ei fastsett akseptgrense. Både haldne og uhaldne aksar kan liggje i det.

Eit omsyn kan vere verna gjennom eit anna (6.675). Når vern manglar, kan grunnen vere feil estimat, utelaten prioritering, utilstrekkeleg kapasitet, konflikt eller endra vilkår. Sjølv når alle veit kva form som vil svikte, står risikoen att så lenge den forma har positivt valsannsyn (7.152).
\end{formproposition}

\begin{formproposition}{7.24}{t}
Læra kan seie kva som tel for kven, kor mykje det tel i valet, og kva det fører til. Kor mykje det bør telje, seier ho ikkje. For kvar fordeling av vekter ho skildrar, står spørsmålet ope om fordelinga er rett.

Ei sterk vekt i valmodellen er eit mål på innverknad, ikkje eit bevis på rett. Setningane om omsorg seier kva som må vere tilfelle for at den andre skal kunne telje, og kva som hender når han ikkje gjer det. Dei seier ikkje at han bør.
\end{formproposition}

\begin{formproposition}{7.241}{t}
Omsorga minskar avviket hos den andre berre når ho har vekt hos nokon med innverknad (6.7), eit bilete som ordnar alternativa slik avviket hans gjer (6.711), og eit hald som kan gjennomførast (5.21). Om biletet ordnar rett, kan berre svaret hans eller ei uavhengig prøve av følgjene avgjere (6.712).\kjelder{68}

Utan vekt er omsorga eit ønske. Utan hald er ho avmakt. Med feilvendt bilete er ho skade. Utan veg tilbake veit ingen kva ho er.

Den som ikkje kan svare, har berre prøva av følgjene og den som svarar for han. Den som svarar for ein annan, har eit bilete, ikkje eit svar (6.713).
\end{formproposition}

\begin{formproposition}{7.242}{t}
Vekt i valet er ikkje hald. Vurderinga \(\Phi\) er ein sum, og i ein sum kan små gevinstar hos mange vege opp eit stort avvik hos éin (6.6). Haldet er ikkje ein sum: det gjeld kvart avvik for seg, mot toleransen til den som ber det (5.231).

Om den andre talde, viser vektene. Om han vart halden, viser avviket hans mot toleransen hans. Ei form kan ha låg samla spenning og likevel presse éin over grensa.

Skal summen ikkje kunne vege opp eit slikt avvik, må forma vere stengd ute frå \(D\), ikkje berre vekta ned (6.003).
\end{formproposition}

\begin{formproposition}{7.243}{t}
Avviket hos den andre kan minskast på to måtar: ved å flytte forma mot målet hans, eller ved å flytte målet hans mot forma (5.131, 7.2). Den som flyttar målet, har endra målestokken.

Setningane om omsorg gjeld difor berre når avviket vert målt mot ein målregion forma ikkje sjølv har flytta: målet den andre hadde før forma verka på han, eller det han held fast ved når alternativa står opne for han. Vender målet tilbake når alternativa kjem til, var endringa tilpassing til forma, ikkje eit svar på henne.
\end{formproposition}

\begin{formproposition}{7.25}{t}
Omsorga kan rette seg sjølv gjennom prøver, men kan ikkje gje seg sjølv endeleg stadfesting på vegner av den andre (6.712). Svaret frå den som lever med forma, kan endre både målet, modellen og det som må haldast.

Å lukke sløyfa er å la dette svaret endre arbeidet.
\end{formproposition}

\begin{formproposition}{7.26}{a}
Denne læra set éi forplikting for design som fag: omsorga skal målast mot verknaden, ikkje mot bestillinga. Den som ber vesentlege følgjer av ei form (7.011), skal ha vekt i valet og ein veg tilbake for svaret sitt. Ingen kan reknast bort fordi han manglar makt til å setje krava.

Dette er ei forplikting, ikkje ei følgje av rekninga og ikkje ein hypotese. Ho kan ikkje prøvast. Ho kan haldast eller brytast, og om ho vert halden, kan prøvast (7.241).

Der omsorga ikkje rekk, eller eit omsyn må vike for eit anna (2.13), skal det seiast: kven som valde måla, kven som ber risikoen, kven som kan rette arbeidet, og kven som svarar for dei som ikkje kan svare sjølve.
\end{formproposition}

\begin{formproposition}{7.261}{t}
Forpliktinga legg ikkje eit nytt kriterium til. Ho fastset rekkjevidda: designkriteriet (6.02) skal gjelde i heile verknadskjegla (7.011), ikkje berre langs dei omsyna bestillinga nemner.

Kriteriet er ein definisjon. Rekkjevidda er forpliktinga. Langs dei omsyna der kriteriet er møtt, er arbeidet design. Langs dei andre er det formgjeving av vilkåra til andre (6.021, 6.022).
\end{formproposition}

\begin{formproposition}{7.262}{t}
Dei fleste hinner vert delvis haldne av andre (5.0221). Det som held agenten, høyrer til vilkåra for hans eige hald, anten det ligg i lyskjegla hans eller ikkje (5.222). Å halde det ved like er sjølvhald med lengre horisont (5.115).

Dette er omhug for eige hald, ikkje omsorg (5.2821). Grunnen rekk så langt som avhengigheita. Han femner ikkje dei agenten ikkje er avhengig av, og ikkje dei som kjem etter han. Der mange vert haldne av det same, kan kvar late dei andre halde det.
\end{formproposition}

\begin{formproposition}{7.263}{t}
Ein del høyrer til ein felles agent berre langs dei omsyna der avviket hans kjem tilbake som avvik i det felles (6.12, 6.1213). Der det ikkje kjem tilbake, er han omgjevnad for heilskapen.

Ei ordning som legg avvika sine på dei som ikkje kan svare, går ikkje dermed i oppløysing. Ho har sett dei utanfor og held fram som éin agent utan dei. Skalaspranget flyttar grensa for omsorga. Det fjernar henne ikkje (6.75).

Tapet kjem tilbake berre der ordninga sjølv vert halden av dei ho legg avvika på (7.262), og då med ei forseinking som kan liggje utanfor det ho føreseier (7.14).
\end{formproposition}

\begin{formproposition}{7.264}{t}
Kvar form som verkar på ein annan, har gjeve han ei vekt i valet, om det så er vekta null (6.671, 6.7). Den som formgjev, kan late spørsmålet liggje. Svaret gjev forma likevel.

Kor mykje den andre bør telje, seier ikkje læra. Kor mykje han talde, viser forma og forløpet (7.02).
\end{formproposition}

\begin{formproposition}{7.27}{t}
Kunnskap frå eit prosjekt kan overførast når me kan vise kva relasjonar og vilkår som får verknaden til å gå att. Same geometri er ikkje alltid nødvendig; same namn på funksjonen er ikkje nok.

Det som varer, er kunnskap om kva som ber og kva endringar det toler.
\end{formproposition}

\begin{formproposition}{7.271}{t}
Eit resultat går inn i det kjende saman med prøvevilkåra sine (1.52). Ein fiasko høyrer til det kjende like mykje som ein suksess. Å ta opp att ei prøve under same vilkår er ei \(K\to K\)-rørsle (1.54), ikkje ei gjentaking: ho viser om avgjerda held når ein annan gjer prøva.

Kvar ny form tek utgangspunkt i ein busett posisjon (1.411). Kvart nytt arbeid tek på same vis utgangspunkt i ei avgjerd i det kjende, og må seie kva for ei. Ei avhandling er eit slikt arbeid: ho nemner avgjerda ho prøver, og lèt etter seg ei avgjerd den neste kan prøve.
\end{formproposition}

\begin{formproposition}{7.272}{t}
Kunnskapsstanden høyrer til ein agent (1.5, 5.118). Det éin har prøvd, er avgjort for ein annan først når prøvevilkåra, dataa og stega følgjer med, så prøva kan takast opp att. Utan dei er resultatet for alle andre ein påstand: tenkt, ikkje avgjort.

Det faget kjenner, er det fleire enn éin kan prøve (1.5311).
\end{formproposition}

\begin{formproposition}{7.273}{t}
Ingen ser si eiga utelating innanfrå (1.53). Ho kan peikast ut frå ein annan kunnskapsstand (1.531). Ei etterprøving gjord av dei som deler utelatingane, finn dei ikkje.

Uavhengig er den som prøver utan å ha valt utfallet og utan å bere tapet ved eit motdøme.
\end{formproposition}

\begin{formproposition}{7.274}{t}
To prøver med ulike aksar, toleransar eller grunnfordelingar har ikkje prøvd det same (1.321, 7.151). Eitt resultat kan rette eit anna berre når dei er målte på same måte.

Eit måleinstrument er sjølv ein påstand: at det skil det aksen skil, og gjev same svar når ingenting er endra. Påstanden må prøvast før instrumentet vert brukt til å prøve noko anna (6.6732).\kjelder{168} Inntil instrumenta finst, er føreseiingane i 7.174 oppgåver, ikkje prøver.
\end{formproposition}

\begin{formproposition}{7.275}{t}
Ei prøve avgjer noko berre når utfallet kunne ha vorte eit anna (1.63). Ein verknad er ein skilnad frå eit oppgjeve alternativ (1.13). Utan alternativet viser arbeidet berre at det vart gjort.\kjelder{167}

Ei prøve som ikkje kunne ha gått imot bestillinga, er ikkje ei prøve av forma. Ho er ein del av bestillinga (6.026).
\end{formproposition}

\begin{formproposition}{7.276}{t}
Ei prøve er sjølv ei form: ho endrar vilkåra for dei ho vert gjord på (7.2). Dei ber følgjer av henne. Forskinga står difor under same kriterium som forma (6.02, 7.26).

Ansvaret følgjer verknaden, ikkje prosjekttida (7.01).
\end{formproposition}

\begin{formproposition}{7.277}{t}
Definisjonane er uavhengige av kva agenten er laga av (5.111). Det læra finn om steg, prioritering, hald og omsorg, gjeld difor kvar vilkåra er oppfylte: i vev, verkstader, føretak og maskiner. Andre fag kan prøve det på sitt eige materiale, og eit motdøme derfrå tel like mykje.

Eit funn som berre held der ordet design vert brukt, er eit funn om ordet.
\end{formproposition}

\begin{formproposition}{7.28}{t}
Eit fag er éin agent langs eit omsyn berre der svikten hos éin kjem tilbake som avvik hos dei andre (6.12, 6.1213). For kunnskapen er etterprøvinga denne tilbakeføringa. For følgjene er ansvaret det: ei ordning som fører avviket hos den som ber følgja, tilbake som avvik hos den som formgav.

Ansvaret lagar den avhengigheita som manglar (7.262). Utan etterprøving og ansvar er faget mange agentar under eitt namn.
\end{formproposition}

\begin{formproposition}{7.3}{t}
Desse proposisjonane er sjølve ei form. Dei ordnar skilnader, og dei viser kor langt omhugen rakk. Dei må tole den prøva dei krev av faget.
\end{formproposition}

\begin{formproposition}{7.301}{t}
Eit fundament for design skal bere meir kunnskap enn det sjølv rommar. Når nye former og livsvilkår kjem til, skal det gje vegar til å forstå dei og til å rette det som ikkje held.
\end{formproposition}

\begin{formproposition}{7.31}{t}
Teksten kan ha oversett samanhengar som avgjer om stega ber. Ein påvist feil må rettast der slutninga skjer; han kan ikkje gjerast uskuldig ved å kallast endå eit døme på endeleg omhug.
\end{formproposition}

\begin{formproposition}{7.32}{t}
Ein lesar kan finne eit motdøme, eit betre mellomledd eller ei prøve som skil forklaringar teksten lèt falle saman. Lesinga er eit svar i sløyfa når ho endrar kva læra kan forklare og halde.
\end{formproposition}

\begin{formproposition}{7.33}{i}
Når eit trinn sviktar under foten, hjelper det lite å gjere stigen lengre.
\end{formproposition}

\begin{formproposition}{7.34}{t}
Fjellet er ikkje i kartet. Men den som les kartet, står på fjellet.
\end{formproposition}

\begin{formproposition}{7.4}{t}
Forma stoggar. Verknaden held fram.
\end{formproposition}
